\documentclass[aspectratio=169]{beamer}\usepackage[]{graphicx}\usepackage[]{xcolor}
% maxwidth is the original width if it is less than linewidth
% otherwise use linewidth (to make sure the graphics do not exceed the margin)
\makeatletter
\def\maxwidth{ %
  \ifdim\Gin@nat@width>\linewidth
    \linewidth
  \else
    \Gin@nat@width
  \fi
}
\makeatother

\definecolor{fgcolor}{rgb}{0.345, 0.345, 0.345}
\newcommand{\hlnum}[1]{\textcolor[rgb]{0.686,0.059,0.569}{#1}}%
\newcommand{\hlsng}[1]{\textcolor[rgb]{0.192,0.494,0.8}{#1}}%
\newcommand{\hlcom}[1]{\textcolor[rgb]{0.678,0.584,0.686}{\textit{#1}}}%
\newcommand{\hlopt}[1]{\textcolor[rgb]{0,0,0}{#1}}%
\newcommand{\hldef}[1]{\textcolor[rgb]{0.345,0.345,0.345}{#1}}%
\newcommand{\hlkwa}[1]{\textcolor[rgb]{0.161,0.373,0.58}{\textbf{#1}}}%
\newcommand{\hlkwb}[1]{\textcolor[rgb]{0.69,0.353,0.396}{#1}}%
\newcommand{\hlkwc}[1]{\textcolor[rgb]{0.333,0.667,0.333}{#1}}%
\newcommand{\hlkwd}[1]{\textcolor[rgb]{0.737,0.353,0.396}{\textbf{#1}}}%
\let\hlipl\hlkwb

\usepackage{framed}
\makeatletter
\newenvironment{kframe}{%
 \def\at@end@of@kframe{}%
 \ifinner\ifhmode%
  \def\at@end@of@kframe{\end{minipage}}%
  \begin{minipage}{\columnwidth}%
 \fi\fi%
 \def\FrameCommand##1{\hskip\@totalleftmargin \hskip-\fboxsep
 \colorbox{shadecolor}{##1}\hskip-\fboxsep
     % There is no \\@totalrightmargin, so:
     \hskip-\linewidth \hskip-\@totalleftmargin \hskip\columnwidth}%
 \MakeFramed {\advance\hsize-\width
   \@totalleftmargin\z@ \linewidth\hsize
   \@setminipage}}%
 {\par\unskip\endMakeFramed%
 \at@end@of@kframe}
\makeatother

\definecolor{shadecolor}{rgb}{.97, .97, .97}
\definecolor{messagecolor}{rgb}{0, 0, 0}
\definecolor{warningcolor}{rgb}{1, 0, 1}
\definecolor{errorcolor}{rgb}{1, 0, 0}
\newenvironment{knitrout}{}{} % an empty environment to be redefined in TeX

\usepackage{alltt}

% Set lecture number for later use


% Part common to all the lectures
\subtitle{MATH 2740 -- Mathematics of Data Science -- Unit U05}
\author{\texorpdfstring{Julien Arino\newline\url{julien.arino@umanitoba.ca}}{Julien Arino}}
\institute{Department of Mathematics @ University of Manitoba}

% Title of the lecture
\title{Principal component analysis}



\usetheme{default}
% Slide setup, colour independent

\usepackage{amsmath,amssymb,amsthm}
\usepackage[utf8]{inputenc}
\usepackage{colortbl}
\usepackage{bm}
\usepackage{xcolor}
\usepackage{dsfont}
\usepackage{setspace}
% To use \ding{234} and the like
\usepackage{pifont}
% To cross reference between slide files
\usepackage{zref-xr,zref-user}
% Use something like
% \zexternaldocument{fileI}
% in the tex files. And cite using \zref instead of \ref

% Cross-reference system - see CROSS-REFERENCE-SETUP.md for manual setup instructions
\usepackage{booktabs}
\usepackage{marvosym}
\usepackage{cancel}
%\usepackage{transparent}
% Make doi clickable in the bibliography?
\usepackage{doi}

\usepackage[T1]{fontenc}

\usepackage{longtable}

% For heavier titles
\usepackage{helvet} % Enables Helvetica font family


% Fields and the like
\def\IC{\mathbb{C}}
\def\IE{\mathbb{E}}
\def\IF{\mathbb{F}}
\def\II{\mathbb{I}}
\def\IJ{\mathbb{J}}
\def\IK{\mathbb{K}}
\def\IM{\mathbb{M}}
\def\IN{\mathbb{N}}
\def\IP{\mathbb{P}}
\def\IR{\mathbb{R}}
\newcommand{\IRplus}{\mathbb{R}_{\ge 0}}
\def\IZ{\mathbb{Z}}
\def\11{\mathds{1}}


% Bold lowercase
\def\ba{\bm{a}}
\def\bb{\bm{b}}
\def\bc{\bm{c}}
\def\bd{\bm{d}}
\def\be{\bm{e}}
\def\bf{\bm{f}}
\def\bg{\bm{g}}
\def\bh{\bm{h}}
\def\bi{\bm{i}}
\def\bj{\bm{j}}
\def\bk{\bm{k}}
\def\bn{\bm{n}}
\def\bp{\bm{p}}
\def\br{\bm{r}}
\def\bs{\bm{s}}
\def\bu{\bm{u}}
\def\bv{\bm{v}}
\def\bw{\bm{w}}
\def\bx{\bm{x}}
\def\by{\bm{y}}
\def\bz{\bm{z}}
\newcommand{\vect}[1]{\bm{#1}}

% Bold capitals
\def\bB{\bm{B}}
\def\bD{\bm{D}}
\def\bE{\bm{E}}
\def\bF{\bm{F}}
\def\bG{\bm{G}}
\def\bI{\bm{I}}
\def\bL{\bm{L}}
\def\bN{\bm{N}}
\def\bP{\bm{P}}
\def\bR{\bm{R}}
\def\bS{\bm{S}}
\def\bT{\bm{T}}
\def\bX{\bm{X}}

% Bold numbers
\def\b0{\bm{0}}

% Bold greek
\bmdefine{\bmu}{\bm{\mu}}
\def\bphi{\bm{\phi}}
\def\bvarphi{\bm{\varphi}}
\def\bPi{\bm{\Pi}}
\def\bGamma{\bm{\Gamma}}

% Bold red sentence
\def\boldred#1{{\color{red}\textbf{#1}}}
\def\defword#1{{\color{orange}\textbf{#1}}}

% Caligraphic letters
\def\A{\mathcal{A}}
\def\B{\mathcal{B}}
\def\C{\mathcal{C}}
\def\D{\mathcal{D}}
\def\E{\mathcal{E}}
\def\F{\mathcal{F}}
\def\G{\mathcal{G}}
\def\H{\mathcal{H}}
\def\I{\mathcal{I}}
\def\L{\mathcal{L}}
\def\M{\mathcal{M}}
\def\N{\mathcal{N}}
\def\P{\mathcal{P}}
\def\R{\mathcal{R}}
\def\S{\mathcal{S}}
\def\T{\mathcal{T}}
\def\U{\mathcal{U}}
\def\V{\mathcal{V}}

% Adding space for prime (') where needed
\def\pprime{\,'}
% Adding space for star (\star) where needed
\def\pstar{{\,\star}}

% tt font for code
\def\code#1{{\tt #1}}

% i.e., e.g.
\def\eg{\emph{e.g.}}
\def\ie{\emph{i.e.}}


% Operators and special symbols
\def\nbOne{{\mathchoice {\rm 1\mskip-4mu l} {\rm 1\mskip-4mu l}
{\rm 1\mskip-4.5mu l} {\rm 1\mskip-5mu l}}}
\def\cov{\ensuremath{\mathsf{cov}}}
\def\Var{\ensuremath{\mathsf{Var}\ }}
\def\Im{\textrm{Im}\;}
\def\Re{\textrm{Re}\;}
\def\det{\ensuremath{\mathsf{det}}}
\def\diag{\ensuremath{\mathsf{diag}}}
\def\nullspace{\ensuremath{\mathsf{null}}}
\def\nullity{\ensuremath{\mathsf{nullity}}}
\def\rank{\ensuremath{\mathsf{rank}}}
\def\range{\ensuremath{\mathsf{range}}}
\def\sgn{\ensuremath{\mathsf{sgn}}}
\def\Span{\ensuremath{\mathsf{span}}}
\def\tr{\ensuremath{\mathsf{tr}}}
\def\imply{$\Rightarrow$}
\def\restrictTo#1#2{\left.#1\right|_{#2}}
\newcommand{\parallelsum}{\mathbin{\!/\mkern-5mu/\!}}
\def\dsum{\mathop{\displaystyle \sum }}%
\def\dind#1#2{_{\substack{#1\\ #2}}}

\newcommand{\Qmatrix}[1]{%
  \begin{pmatrix}#1\end{pmatrix}%
}

\DeclareMathOperator{\GL}{GL}
\DeclareMathOperator{\Rel}{Re}
\def\Nt#1{\left|\!\left|\!\left|#1\right|\!\right|\!\right|}
\newcommand{\tripbar}{|\! |\! |}



% The beamer bullet (in base colour)
\def\bbullet{\leavevmode\usebeamertemplate{itemize item}\ }

% Theorems and the like
\newtheorem{importantdefinition}[theorem]{Definition}
\newtheorem{proposition}[theorem]{Proposition}
\newtheorem{property}[theorem]{Property}
\newtheorem{importantproperty}[theorem]{Property}
\newtheorem{importanttheorem}[theorem]{Theorem}
%\newtheorem{lemma}[theorem]{Lemma}
%\newtheorem{corollary}[theorem]{Corollary}
\newtheorem{remark}[theorem]{Remark}
\setbeamertemplate{theorems}[numbered]
%\setbeamertemplate{theorems}[ams style]

%
%\usecolortheme{orchid}
%\usecolortheme{orchid}

\def\red{\color[rgb]{1,0,0}}
\def\blue{\color[rgb]{0,0,1}}
\def\green{\color[rgb]{0,1,0}}

% Fix skipping lines after items in the bibliography
\setbeamertemplate{bibliography entry title}{}
\setbeamertemplate{bibliography entry location}{}
\setbeamertemplate{bibliography entry note}{}

% Get rid of navigation stuff
\setbeamertemplate{navigation symbols}{}

% Set footline/header line
\setbeamertemplate{footline}
{%
\quad p. \insertpagenumber \quad--\quad \insertsection\vskip2pt
}
% \setbeamertemplate{headline}
% {%
% \quad\insertsection\hfill p. \insertpagenumber\quad\mbox{}\vskip2pt
% }


\makeatletter
\newlength\beamerleftmargin
\setlength\beamerleftmargin{\Gm@lmargin}
\makeatother

% Colours for special pages
\def\extraContent{yellow!20}


%%%%%%%%%%%%%%%%%
\usepackage{tikz}
\usetikzlibrary{shapes,arrows}
\usetikzlibrary{positioning}
\usetikzlibrary{shapes.symbols,shapes.callouts,patterns}
\usetikzlibrary{calc,fit}
\usetikzlibrary{backgrounds}
\usetikzlibrary{decorations.pathmorphing,fit,petri}
\usetikzlibrary{automata}
\usetikzlibrary{fadings}
\usetikzlibrary{patterns,hobby}
\usetikzlibrary{backgrounds,fit,petri}
\usetikzlibrary{tikzmark}

\usepackage{pgfplots}
\pgfplotsset{compat=1.6}
\pgfplotsset{ticks=none}

\usetikzlibrary{decorations.markings}
\usetikzlibrary{arrows.meta}
\tikzset{>=stealth}

% For tikz
\tikzstyle{cloud} = [draw, ellipse,fill=red!20, node distance=0.87cm,
minimum height=2em]
\tikzstyle{line} = [draw, -latex']


%%% For max frame images
\newenvironment{changemargin}[2]{%
\begin{list}{}{%
\setlength{\topsep}{0pt}%
\setlength{\leftmargin}{#1}%
\setlength{\rightmargin}{#2}%
\setlength{\listparindent}{\parindent}%
\setlength{\itemindent}{\parindent}%
\setlength{\parsep}{\parskip}%
}%
\item[]}{\end{list}}


% Make one image take up the entire slide content area in beamer,.:
% centered/centred full-screen image, with title:
% This uses the whole screen except for the 1cm border around it
% all. 128x96mm
\newcommand{\titledFrameImage}[2]{
\begin{frame}{#1}
%\begin{changemargin}{-1cm}{-1cm}
\begin{center}
\includegraphics[width=108mm,height=\textheight,keepaspectratio]{#2}
\end{center}
%\end{changemargin}
\end{frame}
}

% Make one image take up the entire slide content area in beamer.:
% centered/centred full-screen image, no title:
% This uses the whole screen except for the 1cm border around it
% all. 128x96mm
\newcommand{\plainFrameImage}[1]{
\begin{frame}[plain]
%\begin{changemargin}{-1cm}{-1cm}
\begin{center}
\includegraphics[width=108mm,height=76mm,keepaspectratio]{#1}
\end{center}
%\end{changemargin}
\end{frame}
}

% Make one image take up the entire slide area, including borders, in beamer.:
% centered/centred full-screen image, no title:
% This uses the entire whole screen
\newcommand{\maxFrameImage}[1]{
\begin{frame}[plain]
\begin{changemargin}{-1cm}{-1cm}
\begin{center}
\includegraphics[width=\paperwidth,height=\paperheight,keepaspectratio]
{#1}
\end{center}
\end{changemargin}
\end{frame}
}

% This uses the entire whole screen (to include in frame)
\newcommand{\maxFrameImageNoFrame}[1]{
\begin{changemargin}{-1cm}{-1cm}
\begin{center}
\includegraphics[width=\paperwidth,height=0.99\paperheight,keepaspectratio]
{#1}
\end{center}
\end{changemargin}
}

% Make one image take up the entire slide area, including borders, in beamer.:
% centered/centred full-screen image, no title:
% This uses the entire whole screen
\newcommand{\maxFrameImageColor}[2]{
\begin{frame}[plain]
\setbeamercolor{normal text}{bg=#2!20}
\begin{changemargin}{-1cm}{-1cm}
\begin{center}
\includegraphics[width=\paperwidth,height=\paperheight,keepaspectratio]
{#1}
\end{center}
\end{changemargin}
\end{frame}
}


\usepackage{tikz}
\usetikzlibrary{patterns,hobby,matrix}
\usepackage{pgfplots}
\pgfplotsset{compat=1.6}
\pgfplotsset{ticks=none}

\usetikzlibrary{backgrounds}
\usetikzlibrary{decorations.markings}
\usetikzlibrary{arrows.meta}
\tikzset{>=stealth}

\tikzset{
  clockwise arrows/.style={
    postaction={
      decorate,
      decoration={
        markings,
        mark=between positions 0.1 and 0.9 step 40pt with {\arrow{>}},
   }}}}


% New integrated section command: creates section and section slide
\newcommand{\Ssection}[2]{
\section{#1}
\begin{frame}[noframenumbering,plain]
  \begin{tikzpicture}[remember picture,overlay]
    \node[above right,inner sep=0pt,opacity=0.2] at (current page.south west)
    {
        \includegraphics[height=\paperheight,width=\paperwidth]{#2}
    };
  \end{tikzpicture}
  \setbeamercolor{section in toc}{fg=section_page_list_colour}
  \setbeamerfont{section in toc}{size=\Large,series=\bfseries}
  \setbeamertemplate{section in toc shaded}[default][60]
  \tableofcontents[
    currentsection,
    sectionstyle=show/shaded,
    subsectionstyle=show/hide/hide,
    subsubsectionstyle=hide/hide/hide]
\end{frame}
\addtocounter{page}{-1}
}

% New integrated section command with subsections: creates section and section slide showing subsections
\newcommand{\SsectionWithSubs}[2]{
\section{#1}
\begin{frame}[noframenumbering,plain]
  \begin{tikzpicture}[remember picture,overlay]
    \node[above right,inner sep=0pt,opacity=0.2] at (current page.south west)
    {
        \includegraphics[height=\paperheight,width=\paperwidth]{#2}
    };
  \end{tikzpicture}
  \setbeamercolor{section in toc}{fg=section_page_list_colour}
  \setbeamerfont{section in toc}{size=\Large,series=\bfseries}
  \setbeamertemplate{section in toc shaded}[default][60]
  \tableofcontents[
    currentsection,
    sectionstyle=show/hide,
    subsectionstyle=show/show/hide,
    subsubsectionstyle=hide/hide/hide]
\end{frame}
\addtocounter{page}{-1}
}

% New integrated subsection command: creates subsection and subsection slide
\newcommand{\Ssubsection}[2]{
\subsection{#1}
\begin{frame}[noframenumbering,plain]
  \begin{tikzpicture}[remember picture,overlay]
    \node[above right,inner sep=0pt,opacity=0.2] at (current page.south west)
    {
        \includegraphics[height=\paperheight,width=\paperwidth]{#2}
    };
  \end{tikzpicture}
  \setbeamercolor{section in toc}{fg=subsection_page_list_colour}
  \setbeamerfont{section in toc}{size=\Large,series=\bfseries}
  \setbeamertemplate{section in toc shaded}[default][60]
  \setbeamerfont{subsection in toc}{series=\bfseries}
  \setbeamertemplate{subsection in toc shaded}[default][50]
  \tableofcontents[
    currentsection,
    sectionstyle=show/hide,
    subsectionstyle=show/shaded/hide,
    subsubsectionstyle=hide/hide/hide]
\end{frame}
\addtocounter{page}{-1}
}

% New integrated subsubsection command: creates subsubsection and subsubsection slide
\newcommand{\Ssubsubsection}[2]{
\subsubsection{#1}
\begin{frame}[noframenumbering,plain]
  \begin{tikzpicture}[remember picture,overlay]
    \node[above right,inner sep=0pt,opacity=0.2] at (current page.south west)
    {
        \includegraphics[height=\paperheight,width=\paperwidth]{#2}
    };
  \end{tikzpicture}
  \setbeamercolor{section in toc}{fg=subsub_header_section}
  \setbeamerfont{section in toc}{size=\Large,series=\bfseries}
  \setbeamertemplate{section in toc shaded}[default][60]
  \setbeamerfont{subsection in toc}{series=\bfseries}
  \setbeamertemplate{subsection in toc shaded}[default][50]
  \setbeamertemplate{subsubsection in toc shaded}[default][50]
  \tableofcontents[
    currentsection,
    sectionstyle=show/hide,
    subsectionstyle=show/hide/hide,
    subsubsectionstyle=show/shaded/hide]
\end{frame}
\addtocounter{page}{-1}
}

% Legacy commands (kept for backward compatibility)
% Beginning of a section
\newcommand{\newSectionSlide}[1]{
\begin{frame}[noframenumbering,plain]
  \begin{tikzpicture}[remember picture,overlay]
    \node[above right,inner sep=0pt,opacity=0.2] at (current page.south west)
    {
        \includegraphics[height=\paperheight,width=\paperwidth]{#1}
    };
  \end{tikzpicture}
  \setbeamercolor{section in toc}{fg=section_page_list_colour}
  \setbeamerfont{section in toc}{size=\Large,series=\bfseries}
  \setbeamertemplate{section in toc shaded}[default][60]
  \tableofcontents[
    currentsection,
    sectionstyle=show/shaded,
    subsectionstyle=show/hide/hide,
    subsubsectionstyle=hide/hide/hide]
\end{frame}
\addtocounter{page}{-1}
}

% Beginning of a section in which we also show subsections
\newcommand{\newSectionWithSubsSlide}[1]{
	\begin{frame}[noframenumbering,plain]
		\begin{tikzpicture}[remember picture,overlay]
			\node[above right,inner sep=0pt,opacity=0.2] at (current page.south west)
			{
				\includegraphics[height=\paperheight,width=\paperwidth]{#1}
			};
		\end{tikzpicture}
		\setbeamercolor{section in toc}{fg=section_page_list_colour}
		\setbeamerfont{section in toc}{size=\Large,series=\bfseries}
		\setbeamertemplate{section in toc shaded}[default][60]
		\tableofcontents[
		currentsection,
		sectionstyle=show/hide,
		subsectionstyle=show/show/hide,
		subsubsectionstyle=hide/hide/hide]
	\end{frame}
	\addtocounter{page}{-1}
}

% Beginning of a subsection
\newcommand{\newSubSectionSlide}[1]{
\begin{frame}[noframenumbering,plain]
  \begin{tikzpicture}[remember picture,overlay]
    \node[above right,inner sep=0pt,opacity=0.2] at (current page.south west)
    {
        \includegraphics[height=\paperheight,width=\paperwidth]{#1}
    };
  \end{tikzpicture}
  \setbeamercolor{section in toc}{fg=subsection_page_list_colour}
  \setbeamerfont{section in toc}{size=\Large,series=\bfseries}
  \setbeamertemplate{section in toc shaded}[default][60]
  \setbeamerfont{subsection in toc}{series=\bfseries}
  \setbeamertemplate{subsection in toc shaded}[default][50]
  \tableofcontents[
    currentsection,
    sectionstyle=show/hide,
    subsectionstyle=show/shaded/hide,
    subsubsectionstyle=hide/hide/hide]
\end{frame}
\addtocounter{page}{-1}
}

% Beginning of a subsubsection
\newcommand{\newSubSubSectionSlide}[1]{
\begin{frame}[noframenumbering,plain]
  \begin{tikzpicture}[remember picture,overlay]
    \node[above right,inner sep=0pt,opacity=0.2] at (current page.south west)
    {
        \includegraphics[height=\paperheight,width=\paperwidth]{#1}
    };
  \end{tikzpicture}
  \setbeamercolor{section in toc}{fg=subsub_header_section}
  \setbeamerfont{section in toc}{size=\Large,series=\bfseries}
  \setbeamertemplate{section in toc shaded}[default][60]
  \setbeamerfont{subsection in toc}{series=\bfseries}
  \setbeamertemplate{subsection in toc shaded}[default][50]
  \setbeamertemplate{subsubsection in toc shaded}[default][50]
  \tableofcontents[
    currentsection,
    sectionstyle=show/hide,
    subsectionstyle=show/hide/hide,
    subsubsectionstyle=show/shaded/hide]
\end{frame}
\addtocounter{page}{-1}
}


   %%%%%%%%%%%
% To have links to parts in the outline
\makeatletter
\AtBeginPart{%
  \addtocontents{toc}{\protect\beamer@partintoc{\the\c@part}{\beamer@partnameshort}{\the\c@page}}%
}
%% number, shortname, page.
\providecommand\beamer@partintoc[3]{%
  \ifnum\c@tocdepth=-1\relax
    % requesting onlyparts.
    \makebox[6em]{Part #1:} \textcolor{green!30!blue}{\hyperlink{#2}{#2}}
    \par
  \fi
}
\define@key{beamertoc}{onlyparts}[]{%
  \c@tocdepth=-1\relax
}
\makeatother%

\newcommand{\nameofthepart}{}
\newcommand{\nupart}[1]%
    {   \part{#1}%
        \renewcommand{\nameofthepart}{#1}%
        {
          \setbeamercolor{background canvas}{bg=orange!50}
          \begin{frame}{#1}%\partpage 
          \hypertarget{\nameofthepart}{}\tableofcontents%
          \end{frame}
        }
    }

% This command creates a title page using TikZ only
\newcommand{\tikztitlepage}[1]{%
\begin{frame}[plain,noframenumbering]
  \begin{tikzpicture}[remember picture,overlay]
    % Background image
    \node[above right,inner sep=0pt,opacity=0.1] 
      at (current page.south west) 
      {\includegraphics[width=\paperwidth,height=\paperheight]{#1}};

    % University logo
    \node[anchor=north east, inner sep=5pt, opacity=0.9] 
      at (current page.north east)
      {\includegraphics[width=0.2\textwidth]{FIGS-slides-admin/UM-logo-horizontal-CMYK.png}};
    
    % Title
    \node[anchor=center, align=center, 
          font=\fontsize{13}{15}\bfseries\color{UMbrown}, 
          text width=0.9\textwidth] 
          at ([yshift=2cm]current page.center)
          {\inserttitle};

      % Authors
      \node[anchor=center, align=center,
        font=\fontsize{10}{12}\bfseries\color{UMbrown},
        text width=0.7\textwidth]
        at ([yshift=0.8cm]current page.center)
        {\insertauthor};

      % Affiliation
      \node[anchor=north, align=center,
        font=\fontsize{9}{11}\color{UMbrown},
        text width=0.7\textwidth]
        at ([yshift=-0.2cm]current page.center)
        {\insertaffiliation};      
    % Date
    \node[anchor=north, align=center, 
          font=\fontsize{12}{16}\bfseries\color{UMbrown},
          text width=0.7\textwidth] 
          at ([yshift=0.2cm]current page.center)
          {\insertdate};

    % Land acknowledgement
    \node[anchor=south, align=justify, 
          font=\footnotesize, text=black, 
          text width=1.1\textwidth] 
          at ([yshift=0.5cm]current page.south)
          {The University of Manitoba campuses are located on original lands of Anishinaabeg, Ininew, Anisininew, Dakota and Dene peoples, and on the National Homeland of the Red River Métis.\\
          We respect the Treaties that were made on these territories, we acknowledge the harms and mistakes of the past, and we dedicate ourselves to move forward in partnership with Indigenous communities in a spirit of Reconciliation and collaboration.};
  \end{tikzpicture}
  \addtocounter{page}{-1}
\end{frame}
}
% The title page with figure
% \newcommand{\titlepagewithfigure}[1]{%
%   \begin{frame}[noframenumbering,plain]
%     \begin{tikzpicture}[remember picture,overlay]
%       \node[above right,inner sep=0pt,opacity=0.1] at (current page.south west)
%       {
%           \includegraphics[height=\paperheight,width=\paperwidth]{#1}
%       };
%       \node[anchor=north east,
%       inner sep=5pt,
%       opacity=0.9] at (current page.north east)
%       {
%           \includegraphics[width=0.2\textwidth]{FIGS-slides-admin/UM-logo-horizontal-CMYK.png}
%       };
%       \node[anchor=south, 
%       align=justify, 
%       text=black, 
%       text width=1.1\textwidth,
%       font=\footnotesize]  (land_acknowledgement)
%       at (current page.south) 
%       {The University of Manitoba campuses are located on original lands of Anishinaabeg, Ininew, Anisininew, Dakota and Dene peoples, and on the National Homeland of the Red River Métis.
%       We respect the Treaties that were made on these territories, we acknowledge the harms and mistakes of the past, and we dedicate ourselves to move forward in partnership with Indigenous communities in a spirit of Reconciliation and collaboration.};  
%       % \node[align=center, anchor=south,
%       % above=0.5cm of land_acknowledgement,
%       % text=black,
%       % font=\bfseries] {\insertdate};
%   \end{tikzpicture}
%   \setbeamercolor{title}{fg=title_page_title_colour}
%   \setbeamerfont{title}{size=\Large,series=\bfseries}
%   \setbeamercolor{author}{fg=title_page_author_colour}
%   \setbeamerfont{author}{size=\large,series=\bfseries}
%   \setbeamercolor{institute}{fg=title_page_institute_colour}
%   \setbeamerfont{institute}{size=\large,series=\bfseries}
%   \setbeamercolor{date}{fg=title_page_date_colour}
%   \setbeamerfont{date}{series=\bfseries}
% 	\titlepage
% \end{frame}
% \addtocounter{page}{-1}
% }

\newcommand{\titlepagewithfigure}[1]{%
  \begin{frame}[noframenumbering,plain]
    \begin{tikzpicture}[remember picture,overlay]
      \node[above right,inner sep=0pt,opacity=0.1] at (current page.south west)
      {
          \includegraphics[height=\paperheight,width=\paperwidth]{#1}
      };
      \node[anchor=north east,
      inner sep=5pt,
      opacity=0.9] at (current page.north east)
      {
          \includegraphics[width=0.2\textwidth]{FIGS-slides-admin/UM-logo-horizontal-CMYK.png}
      };
      \node[anchor=south, 
      align=justify, 
      text=black, 
      text width=1.1\textwidth,
      font=\footnotesize]  (land_acknowledgement)
      at (current page.south) 
      {The University of Manitoba campuses are located on original lands of Anishinaabeg, Ininew, Anisininew, Dakota and Dene peoples, and on the National Homeland of the Red River Métis.
      We respect the Treaties that were made on these territories, we acknowledge the harms and mistakes of the past, and we dedicate ourselves to move forward in partnership with Indigenous communities in a spirit of Reconciliation and collaboration.};  
      % \node[align=center, anchor=south,
      % above=0.5cm of land_acknowledgement,
      % text=black,
      % font=\bfseries] {\insertdate};
  \end{tikzpicture}
  \setbeamercolor{title}{fg=title_page_title_colour}
  \setbeamerfont{title}{size=\Large,series=\bfseries,family=\usefont{T1}{phv}{b}{n}}
  \setbeamercolor{author}{fg=title_page_author_colour}
  \setbeamerfont{author}{size=\large,series=\bfseries,family=\usefont{T1}{phv}{b}{n}}
  \setbeamercolor{institute}{fg=title_page_institute_colour}
  \setbeamerfont{institute}{size=\large,series=\bfseries,family=\usefont{T1}{phv}{b}{n}}
  \setbeamercolor{date}{fg=title_page_date_colour}
  \setbeamerfont{date}{series=\bfseries,family=\usefont{T1}{phv}{b}{n}}
	\titlepage
\end{frame}
\addtocounter{page}{-1}
}
% The outline page, with figure
% \newcommand{\outlinepage}[1]{%
% \begin{frame}[noframenumbering,plain]
%   \begin{tikzpicture}[remember picture,overlay]
%     \node[above right,inner sep=0pt,opacity=0.2] at (current page.south west)
%     {
%         \includegraphics[height=\paperheight,width=\paperwidth]{#1}
%     };
%   \end{tikzpicture}
%   \setbeamercolor{section in toc}{fg=outline_page_list_colour}
%   \setbeamerfont{section in toc}{size=\Large,series=\bfseries,family=\sffamily}
%   \frametitle{\textcolor{outline_page_title_colour}{\LARGE\bfseries Outline}}
%   \tableofcontents[hideallsubsections]
% \end{frame}
% \addtocounter{page}{-1}
% }
% The outline page, with figure
\newcommand{\outlinepage}[1]{%
\begin{frame}[noframenumbering,plain]
  \begin{tikzpicture}[remember picture,overlay]
    \node[above right,inner sep=0pt,opacity=0.2] at (current page.south west)
    {
        \includegraphics[height=\paperheight,width=\paperwidth]{#1}
    };
  \end{tikzpicture}
  \setbeamercolor{section in toc}{fg=outline_page_list_colour}
  % Use Helvetica Bold only for the outline slide TOC
  \setbeamerfont{section in toc}{size=\Large,family=\usefont{T1}{phv}{b}{n}}
  % Use Helvetica Bold for the outline title
  \frametitle{\textcolor{outline_page_title_colour}{\usefont{T1}{phv}{b}{n}\LARGE Outline}}
  \tableofcontents[hideallsubsections]
\end{frame}
\addtocounter{page}{-1}
}


%\let\oldsection\section
%\renewcommand{\section}[2]{\oldsection[#1]\newSectionSlide[#2]}


%%%%%%%%%%%%%%%%%%%%%
% CUSTOM SLIDE BACKGROUNDS
%%%%%%%%%%%%%%%%%%%%%
% Define custom background templates for different colors
\defbeamertemplate*{background canvas}{blue}{%
  \color{blue!15}\vrule width\paperwidth height\paperheight%
}
\defbeamertemplate*{background canvas}{green}{%
  \color{green!15}\vrule width\paperwidth height\paperheight%
}
\defbeamertemplate*{background canvas}{red}{%
  \color{red!15}\vrule width\paperwidth height\paperheight%
}
\defbeamertemplate*{background canvas}{yellow}{%
  \color{yellow!20}\vrule width\paperwidth height\paperheight%
}
\defbeamertemplate*{background canvas}{purple}{%
  \color{purple!15}\vrule width\paperwidth height\paperheight%
}
\defbeamertemplate*{background canvas}{orange}{%
  \color{orange!20}\vrule width\paperwidth height\paperheight%
}

% Define keys for the different background options
\makeatletter
\define@key{beamerframe}{blue}[true]{\setbeamertemplate{background canvas}[blue]}
\define@key{beamerframe}{green}[true]{\setbeamertemplate{background canvas}[green]}
\define@key{beamerframe}{red}[true]{\setbeamertemplate{background canvas}[red]}
\define@key{beamerframe}{yellow}[true]{\setbeamertemplate{background canvas}[yellow]}
\define@key{beamerframe}{purple}[true]{\setbeamertemplate{background canvas}[purple]}
\define@key{beamerframe}{orange}[true]{\setbeamertemplate{background canvas}[orange]}
\makeatother

% Reset to normal background for all frames by default
\BeforeBeginEnvironment{frame}{\setbeamertemplate{background canvas}[default]}

\newcommand{\punnett}[2]{
    \begin{center}
    \renewcommand{\arraystretch}{1.5} % Add space to rows
    \begin{tabular}{|c c | c|c|}
        \multicolumn{2}{c}{} & \multicolumn{2}{c}{\textbf{Father}} \\
        \multicolumn{2}{c}{} & #1 \\ \hline
        #2
    \end{tabular}
    \end{center}
}

% Colour definitions for Punnett squares
\usepackage{multirow}
\usepackage{colortbl}
\colorlet{punnettorange}{orange!30}
\colorlet{punnettblack}{black!50}
\colorlet{punnetttortie}{orange!60!black}

\usetikzlibrary{positioning, decorations.pathreplacing, arrows.meta}

% --- TIKZ STYLES FOR CAT DIAGRAMS ---
\tikzset{
    % This is the main placeholder style for the cat images
    catnode/.style={
        draw, 
        rectangle, 
        rounded corners, 
        minimum height=1.5cm, 
        minimum width=2.5cm, 
        align=center, 
        font=\small\bfseries
    },
    % Allele styles (the X and Y at the top/side)
    allele/.style={font=\Large\bfseries},
    female/.style={allele, color=purple!80!black},
    male/.style={allele, color=blue!80!black},
    % --- Placeholder styles for different cats ---
    % YOU CAN REPLACE THE CONTENTS OF THESE NODES WITH YOUR IMAGES
    orange_cat/.style={
        catnode, 
        fill=orange!30, 
        text=black
    },
    black_cat/.style={
        catnode, 
        fill=black!70, 
        text=white
    },
    tortie_cat/.style={
        catnode, 
        fill=orange!50!black, % A mix for tortoiseshell
        text=white
    }
}

\usetikzlibrary{
    arrows.meta, % For nicer arrow heads (e.g., -Latex)
    positioning, % For relative node placement (e.g., above=of)
    automata     % For state diagrams, loops
}
% --- END TIKZ STYLES ---

% Common date for all slides
\date{Fall 2026}
% Counter for random transition slides generated by R
\newcounter{sssectionbg}
\setcounter{sssectionbg}{1}
\newcommand{\SSsection}[1]{
  \Ssection{#1}{FIGS/random_bg_\thesssectionbg.pdf}
  \stepcounter{sssectionbg}
}

% Random transition slide generated by R for the title page
\newcommand{\SStitlepage}{
  \titlepagewithfigure{FIGS/random_title_bg.pdf}
}

% Random transition slide generated by R for the outline page
\newcommand{\SSoutlinepage}{
  \outlinepage{FIGS/random_outline_bg.pdf}
}
\newcommand{\SSsubsection}[1]{
  \Ssubsection{#1}{FIGS/random_bg_\thesssectionbg.pdf}
  \stepcounter{sssectionbg}
}

\newcommand{\SSsubsubsection}[1]{
  \Ssubsubsection{#1}{FIGS/random_bg_\thesssectionbg.pdf}
  \stepcounter{sssectionbg}
}



\usecolortheme{orchid}
%% Listings
\usepackage{listings}
\definecolor{mygreen}{rgb}{0,0.6,0}
\definecolor{mygray}{rgb}{0.5,0.5,0.5}
\definecolor{mymauve}{rgb}{0.58,0,0.82}
\definecolor{mygold}{rgb}{1,0.843,0}
\definecolor{myblue}{rgb}{0.537,0.812,0.941}

\definecolor{mygold2}{RGB}{120,105,22}
\definecolor{mygrey2}{RGB}{50,50,50}

\definecolor{lgreen}{rgb}{0.6,0.9,.6}
\definecolor{lred}{rgb}{1,0.5,.5}

\lstloadlanguages{R}
\lstset{ %
  language=R,
  backgroundcolor=\color{black!05},   % choose the background color
  basicstyle=\footnotesize\ttfamily,        % size of fonts used for the code
  breaklines=true,                 % automatic line breaking only at whitespace
  captionpos=b,                    % sets the caption-position to bottom
  commentstyle=\color{mygreen},    % comment style
  escapeinside={\%*}{*)},          % if you want to add LaTeX within your code
  keywordstyle=\color{red},       % keyword style
  stringstyle=\color{mygold},     % string literal style
  keepspaces=true,
  columns=fullflexible,
  tabsize=4,
}
% Could also do (in lstset)
% basicstyle==\fontfamily{pcr}\footnotesize
\lstdefinelanguage{Renhanced}%
  {keywords={abbreviate,abline,abs,acos,acosh,action,add1,add,%
      aggregate,alias,Alias,alist,all,anova,any,aov,aperm,append,apply,%
      approx,approxfun,apropos,Arg,args,array,arrows,as,asin,asinh,%
      atan,atan2,atanh,attach,attr,attributes,autoload,autoloader,ave,%
      axis,backsolve,barplot,basename,besselI,besselJ,besselK,besselY,%
      beta,binomial,body,box,boxplot,break,browser,bug,builtins,bxp,by,%
      c,C,call,Call,case,cat,category,cbind,ceiling,character,char,%
      charmatch,check,chol,chol2inv,choose,chull,class,close,cm,codes,%
      coef,coefficients,co,col,colnames,colors,colours,commandArgs,%
      comment,complete,complex,conflicts,Conj,contents,contour,%
      contrasts,contr,control,helmert,contrib,convolve,cooks,coords,%
      distance,coplot,cor,cos,cosh,count,fields,cov,covratio,wt,CRAN,%
      create,crossprod,cummax,cummin,cumprod,cumsum,curve,cut,cycle,D,%
      data,dataentry,date,dbeta,dbinom,dcauchy,dchisq,de,debug,%
      debugger,Defunct,default,delay,delete,deltat,demo,de,density,%
      deparse,dependencies,Deprecated,deriv,description,detach,%
      dev2bitmap,dev,cur,deviance,off,prev,,dexp,df,dfbetas,dffits,%
      dgamma,dgeom,dget,dhyper,diag,diff,digamma,dim,dimnames,dir,%
      dirname,dlnorm,dlogis,dnbinom,dnchisq,dnorm,do,dotplot,double,%
      download,dpois,dput,drop,drop1,dsignrank,dt,dummy,dump,dunif,%
      duplicated,dweibull,dwilcox,dyn,edit,eff,effects,eigen,else,%
      emacs,end,environment,env,erase,eval,equal,evalq,example,exists,%
      exit,exp,expand,expression,External,extract,extractAIC,factor,%
      fail,family,fft,file,filled,find,fitted,fivenum,fix,floor,for,%
      For,formals,format,formatC,formula,Fortran,forwardsolve,frame,%
      frequency,ftable,ftable2table,function,gamma,Gamma,gammaCody,%
      gaussian,gc,gcinfo,gctorture,get,getenv,geterrmessage,getOption,%
      getwd,gl,glm,globalenv,gnome,GNOME,graphics,gray,grep,grey,grid,%
      gsub,hasTsp,hat,heat,help,hist,home,hsv,httpclient,I,identify,if,%
      ifelse,Im,image,\%in\%,index,influence,measures,inherits,install,%
      installed,integer,interaction,interactive,Internal,intersect,%
      inverse,invisible,IQR,is,jitter,kappa,kronecker,labels,lapply,%
      layout,lbeta,lchoose,lcm,legend,length,levels,lgamma,library,%
      licence,license,lines,list,lm,load,local,locator,log,log10,log1p,%
      log2,logical,loglin,lower,lowess,ls,lsfit,lsf,ls,machine,Machine,%
      mad,mahalanobis,make,link,margin,match,Math,matlines,mat,matplot,%
      matpoints,matrix,max,mean,median,memory,menu,merge,methods,min,%
      missing,Mod,mode,model,response,mosaicplot,mtext,mvfft,na,nan,%
      names,omit,nargs,nchar,ncol,NCOL,new,next,NextMethod,nextn,%
      nlevels,nlm,noquote,NotYetImplemented,NotYetUsed,nrow,NROW,null,%
      numeric,\%o\%,objects,offset,old,on,Ops,optim,optimise,optimize,%
      options,or,order,ordered,outer,package,packages,page,pairlist,%
      pairs,palette,panel,par,parent,parse,paste,path,pbeta,pbinom,%
      pcauchy,pchisq,pentagamma,persp,pexp,pf,pgamma,pgeom,phyper,pico,%
      pictex,piechart,Platform,plnorm,plogis,plot,pmatch,pmax,pmin,%
      pnbinom,pnchisq,pnorm,points,poisson,poly,polygon,polyroot,pos,%
      postscript,power,ppoints,ppois,predict,preplot,pretty,Primitive,%
      print,prmatrix,proc,prod,profile,proj,prompt,prop,provide,%
      psignrank,ps,pt,ptukey,punif,pweibull,pwilcox,q,qbeta,qbinom,%
      qcauchy,qchisq,qexp,qf,qgamma,qgeom,qhyper,qlnorm,qlogis,qnbinom,%
      qnchisq,qnorm,qpois,qqline,qqnorm,qqplot,qr,Q,qty,qy,qsignrank,%
      qt,qtukey,quantile,quasi,quit,qunif,quote,qweibull,qwilcox,%
      rainbow,range,rank,rbeta,rbind,rbinom,rcauchy,rchisq,Re,read,csv,%
      csv2,fwf,readline,socket,real,Recall,rect,reformulate,regexpr,%
      relevel,remove,rep,repeat,replace,replications,report,require,%
      resid,residuals,restart,return,rev,rexp,rf,rgamma,rgb,rgeom,R,%
      rhyper,rle,rlnorm,rlogis,rm,rnbinom,RNGkind,rnorm,round,row,%
      rownames,rowsum,rpois,rsignrank,rstandard,rstudent,rt,rug,runif,%
      rweibull,rwilcox,sample,sapply,save,scale,scan,scan,screen,sd,se,%
      search,searchpaths,segments,seq,sequence,setdiff,setequal,set,%
      setwd,show,sign,signif,sin,single,sinh,sink,solve,sort,source,%
      spline,splinefun,split,sqrt,stars,start,stat,stem,step,stop,%
      storage,strstrheight,stripplot,strsplit,structure,strwidth,sub,%
      subset,substitute,substr,substring,sum,summary,sunflowerplot,svd,%
      sweep,switch,symbol,symbols,symnum,sys,status,system,t,table,%
      tabulate,tan,tanh,tapply,tempfile,terms,terrain,tetragamma,text,%
      time,title,topo,trace,traceback,transform,tri,trigamma,trunc,try,%
      ts,tsp,typeof,unclass,undebug,undoc,union,unique,uniroot,unix,%
      unlink,unlist,unname,untrace,update,upper,url,UseMethod,var,%
      variable,vector,Version,vi,warning,warnings,weighted,weights,%
      which,while,window,write,\%x\%,x11,X11,xedit,xemacs,xinch,xor,%
      xpdrows,xy,xyinch,yinch,zapsmall,zip},%
   otherkeywords={!,!=,~,$,*,\%,\&,\%/\%,\%*\%,\%\%,<-,<<-,_,/},%
   alsoother={._$},%
   sensitive,%
   morecomment=[l]\#,%
   morestring=[d]",%
   morestring=[d]'% 2001 Robert Denham
  }%

% Custom color for listings
\definecolor{codegreen}{rgb}{0,0.6,0}
\definecolor{codegray}{rgb}{0.5,0.5,0.5}
\definecolor{codepurple}{rgb}{0.58,0,0.82}
\definecolor{backcolour}{rgb}{0.95,0.95,0.92}

%------------------------------------------------------------
% Listings configuration for Julia
%------------------------------------------------------------
% Define the Julia language for the listings package
\lstdefinelanguage{Julia}
{
  morekeywords={abstract,baremodule,begin,break,catch,const,continue,do,else,elseif,end,export,false,finally,for,function,global,if,import,let,local,macro,module,quote,return,struct,true,try,using,while,in},
  sensitive=true,
  comment=[l]{\#},
  morecomment=[n]{\#=}{=\#}, % For block comments like #= ... =#
  string=[b]",
  morestring=[b]',
  alsoletter={!}, % Allow ! in identifiers for mutating functions
}

% Define the style for the code blocks
\lstdefinestyle{JuliaStyle}{
    backgroundcolor=\color{backcolour},   
    commentstyle=\color{codegreen},
    keywordstyle=\color{magenta},
    numberstyle=\tiny\color{codegray},
    stringstyle=\color{codepurple},
    basicstyle=\ttfamily\footnotesize,
    breakatwhitespace=false,         
    breaklines=true,                 
    captionpos=b,                    
    keepspaces=true,                 
    % numbers=left,                    
    % numbersep=5pt,                  
    showspaces=false,                
    showstringspaces=false,
    showtabs=false,                  
    tabsize=2,
    language=Julia % Use the newly defined Julia language
}
\lstset{style=JuliaStyle}

%%%%%%% 
%% Definitions in yellow boxes
\usepackage{etoolbox}
\setbeamercolor{block title}{use=structure,fg=structure.fg,bg=structure.fg!40!bg}
\setbeamercolor{block body}{parent=normal text,use=block title,bg=block title.bg!20!bg}

% Important definition
\BeforeBeginEnvironment{importantdefinition}{%
	\setbeamercolor{block title}{fg=black,bg=yellow!20!white}
	\setbeamercolor{block body}{fg=black, bg=yellow!05!white}
}
\AfterEndEnvironment{importantdefinition}{
	\setbeamercolor{block title}{use=structure,fg=structure.fg,bg=structure.fg!20!bg}
	\setbeamercolor{block body}{parent=normal text,use=block title,bg=block title.bg!50!bg, fg=black}
}
% Regular (less important) definition
\BeforeBeginEnvironment{definition}{%
	\setbeamercolor{block title}{fg=black,bg=black!15!white}
	\setbeamercolor{block body}{fg=black, bg=black!05!white}
}
\AfterEndEnvironment{definition}{
	\setbeamercolor{block title}{use=structure,fg=structure.fg,bg=structure.fg!20!bg}
	\setbeamercolor{block body}{parent=normal text,use=block title,bg=block title.bg!50!bg, fg=black}
}
% Important theorem (red)
\BeforeBeginEnvironment{importanttheorem}{%
	\setbeamercolor{block title}{fg=black,bg=red!20!white}
	\setbeamercolor{block body}{fg=black, bg=red!05!white}
}
\AfterEndEnvironment{importanttheorem}{
	\setbeamercolor{block title}{use=structure,fg=structure.fg,bg=structure.fg!20!bg}
	\setbeamercolor{block body}{parent=normal text,use=block title,bg=block title.bg!50!bg, fg=black}
}
\BeforeBeginEnvironment{importantproperty}{%
	\setbeamercolor{block title}{fg=black,bg=red!50!white}
	\setbeamercolor{block body}{fg=black, bg=red!30!white}
}
\AfterEndEnvironment{importantproperty}{
	\setbeamercolor{block title}{use=structure,fg=structure.fg,bg=structure.fg!20!bg}
	\setbeamercolor{block body}{parent=normal text,use=block title,bg=block title.bg!50!bg, fg=black}
}

%% Colours for title page, outline and the like
\definecolor{UMbrown}{RGB}{79,44,29}
\definecolor{HighlightGold}{RGB}{230,165,83}
\definecolor{DeepGreen}{RGB}{27,46,28}
\definecolor{MediumGreen}{RGB}{52,128,56}
\definecolor{LightGreen}{RGB}{150,224,154}
\definecolor{AccentRed}{RGB}{196,44,27}
%% These are the titles used in practice
% Title page
\colorlet{title_page_title_colour}{UMbrown}
\colorlet{title_page_author_colour}{UMbrown}
\colorlet{title_page_date_colour}{UMbrown}
\colorlet{title_page_institute_colour}{UMbrown}
% Outline page
\colorlet{outline_page_title_colour}{DeepGreen}
\colorlet{outline_page_list_colour}{UMbrown}
% Section page
\colorlet{section_page_title_colour}{DeepGreen}
\colorlet{section_page_list_colour}{UMbrown}
% Subection page
\colorlet{subsection_page_title_colour}{DeepGreen}
\colorlet{subsection_page_list_colour}{UMbrown}



\IfFileExists{upquote.sty}{\usepackage{upquote}}{}
\begin{document}

% Set up cross-references and counter persistence
\setcounter{theorem}{79}

%%%%%%%%%%%%%%%%%%%%%%%%%%%%%%%%%
%%%%%%%%%%%%%%%%%%%%%%%%%%%%%%%%%
%% TITLE AND OUTLINE
%%%%%%%%%%%%%%%%%%%%%%%%%%%%%%%%%
%%%%%%%%%%%%%%%%%%%%%%%%%%%%%%%%%
\SStitlepage
\SSoutlinepage



\begin{frame}{Dimensionality reduction}
In the previous unit, we saw that SVD could be used to reduce the dimension of data (matrices)
\vfill
Now we look at another dimensionality reduction technique: principal component analysis (PCA)
\vfill
PCA is often used as a blackbox technique, here we take a look at the math behind it
\end{frame}


\begin{frame}{What is PCA?}
Linear algebraic technique 
\vfill
Helps reduce a complex dataset to a lower dimensional one
\vfill
Non-parametric method: does not assume anything about data distribution (distribution from the statistical point of view)
\end{frame}

%%%%%%%%%%%%%%%%%%%%%
%%%%%%%%%%%%%%%%%%%%%
\SSsection{A running example: hockey players}

\begin{frame}[fragile]{A 2D example}
Dataset (\href{https://opendata.stackexchange.com/questions/7793/age-weight-and-height-dataset}{link}) of height and weight of some hockey players
\vfill
\begin{knitrout}
\definecolor{shadecolor}{rgb}{0.969, 0.969, 0.969}\color{fgcolor}\begin{kframe}
\begin{alltt}
\hlcom{# From https://figshare.com/ndownloader/files/5303173}
\hldef{data} \hlkwb{=} \hlkwd{read.csv}\hldef{(}\hlsng{"https://github.com/julien-arino/math-of-data-science/raw/refs/heads/main/DATA/hockey-players.csv"}\hldef{)}
\hlkwd{dim}\hldef{(data)}
\end{alltt}
\begin{verbatim}
## [1] 6292   13
\end{verbatim}
\end{kframe}
\end{knitrout}
\vfill
In case you are wondering, this is a database of ice hockey players at IIHF world championships, 2001-2016, assembled by the dataset's author
\vfill
See some comments \href{https://ikashnitsky.github.io/2017/ice-hockey-players-height/}{here}
\vfill

\end{frame}

\begin{knitrout}
\definecolor{shadecolor}{rgb}{0.969, 0.969, 0.969}\color{fgcolor}\begin{kframe}
\begin{alltt}
\hlkwd{head}\hldef{(data,} \hlkwc{n}\hldef{=}\hlnum{3}\hldef{)}
\end{alltt}
\begin{verbatim}
##   year country no                name position side height weight      birth
## 1 2001     RUS 10    tverdovsky oleg         D    L    185     84 1976-05-18
## 2 2001     RUS  2  vichnevsky vitali         D    L    188     86 1980-03-18
## 3 2001     RUS 26 petrochinin evgeni         D    L    182     95 1976-02-07
##                      club      age cohort      bmi
## 1   anaheim mighty ducks  24.95277   1976 24.54346
## 2   anaheim mighty ducks  21.11978   1980 24.33228
## 3 severstal cherepovetal  25.22930   1976 28.68011
\end{verbatim}
\end{kframe}
\end{knitrout}

\begin{frame}
As usual, it is a good idea to plot this to get a sense of the lay of the land
\end{frame}




\maxFrameImage{FIGS/U05-plot-hockey-1-1.pdf}

\begin{frame}[fragile]
The author of the study is interested in the evolution of weights, so it is likely that the same person will be in the dataset several times
\vfill
Let us check this: first check below will be \code{FALSE} if the number of unique names does not match the number of rows in the dataset

\begin{knitrout}
\definecolor{shadecolor}{rgb}{0.969, 0.969, 0.969}\color{fgcolor}\begin{kframe}
\begin{alltt}
\hlkwd{length}\hldef{(}\hlkwd{unique}\hldef{(data}\hlopt{$}\hldef{name))} \hlopt{==} \hlkwd{dim}\hldef{(data)[}\hlnum{1}\hldef{]}
\end{alltt}
\begin{verbatim}
## [1] FALSE
\end{verbatim}
\begin{alltt}
\hlkwd{length}\hldef{(}\hlkwd{unique}\hldef{(data}\hlopt{$}\hldef{name))}
\end{alltt}
\begin{verbatim}
## [1] 3278
\end{verbatim}
\begin{alltt}
\hlkwd{dim}\hldef{(data)[}\hlnum{1}\hldef{]}
\end{alltt}
\begin{verbatim}
## [1] 6292
\end{verbatim}
\end{kframe}
\end{knitrout}
\end{frame}

\begin{frame}[fragile]
We are not interested in the evolution of weights, so we simplify: if more than one record for someone, take average of recorded weights and heights
\vfill
To be extra careful, could check as well that there are no major variations on player height (homonymies?)
\vfill
\begin{knitrout}
\definecolor{shadecolor}{rgb}{0.969, 0.969, 0.969}\color{fgcolor}\begin{kframe}
\begin{alltt}
\hldef{data_simplified} \hlkwb{=} \hlkwd{data.frame}\hldef{(}\hlkwc{name} \hldef{=} \hlkwd{unique}\hldef{(data}\hlopt{$}\hldef{name))}
\hldef{w} \hlkwb{=} \hlkwd{c}\hldef{()}
\hldef{h} \hlkwb{=} \hlkwd{c}\hldef{()}
\hlkwa{for} \hldef{(n} \hlkwa{in} \hldef{data_simplified}\hlopt{$}\hldef{name) \{}
    \hldef{tmp} \hlkwb{=} \hldef{data[}\hlkwd{which}\hldef{(data}\hlopt{$}\hldef{name} \hlopt{==} \hldef{n),]}
    \hldef{h} \hlkwb{=} \hlkwd{c}\hldef{(h,} \hlkwd{mean}\hldef{(tmp}\hlopt{$}\hldef{height))}
    \hldef{w} \hlkwb{=} \hlkwd{c}\hldef{(w,} \hlkwd{mean}\hldef{(tmp}\hlopt{$}\hldef{weight))}
\hldef{\}}
\hldef{data_simplified}\hlopt{$}\hldef{weight} \hlkwb{=} \hldef{w}
\hldef{data_simplified}\hlopt{$}\hldef{height} \hlkwb{=} \hldef{h}
\end{alltt}
\end{kframe}
\end{knitrout}
\end{frame}

\begin{frame}[fragile]
\begin{knitrout}
\definecolor{shadecolor}{rgb}{0.969, 0.969, 0.969}\color{fgcolor}\begin{kframe}
\begin{alltt}
\hldef{data} \hlkwb{=} \hldef{data_simplified}
\hlkwd{head}\hldef{(data_simplified,} \hlkwc{n} \hldef{=} \hlnum{6}\hldef{)}
\end{alltt}
\begin{verbatim}
##                  name weight height
## 1    tverdovsky oleg    84.0  185.0
## 2  vichnevsky vitali    86.0  188.0
## 3 petrochinin evgeni    95.0  182.0
## 4    zhdan alexander    85.5  178.5
## 5    orekhovsky oleg    88.0  175.0
## 6      zhukov sergei    92.5  193.0
\end{verbatim}
\end{kframe}
\end{knitrout}
\end{frame}



\maxFrameImage{FIGS/U05-plot-hockey-2-1.pdf}

\begin{frame}[fragile]\frametitle{Centre the data}
\begin{knitrout}
\definecolor{shadecolor}{rgb}{0.969, 0.969, 0.969}\color{fgcolor}\begin{kframe}
\begin{alltt}
\hlkwd{mean}\hldef{(data}\hlopt{$}\hldef{weight)}
\end{alltt}
\begin{verbatim}
## [1] 87.71555
\end{verbatim}
\begin{alltt}
\hlkwd{mean}\hldef{(data}\hlopt{$}\hldef{height)}
\end{alltt}
\begin{verbatim}
## [1] 183.8596
\end{verbatim}
\begin{alltt}
\hldef{data}\hlopt{$}\hldef{weight.c} \hlkwb{=} \hldef{data}\hlopt{$}\hldef{weight}\hlopt{-}\hlkwd{mean}\hldef{(data}\hlopt{$}\hldef{weight)}
\hldef{data}\hlopt{$}\hldef{height.c} \hlkwb{=} \hldef{data}\hlopt{$}\hldef{height}\hlopt{-}\hlkwd{mean}\hldef{(data}\hlopt{$}\hldef{height)}
\end{alltt}
\end{kframe}
\end{knitrout}
\end{frame}

\maxFrameImage{FIGS/U05-plot-hockey-centred-1.pdf}



\begin{frame}{Setting things up}
Each player is a row in the data matrix (an \emph{observation}), each feature or variable (\emph{height} and \emph{weight}) is a column
\vfill
After deduplication, we have an 3278$\times 2$ matrix (actually, 3278$\times 4$ if we consider the uncentred and centred features, but we will use one or the other, not both uncentred and centred)
\vfill
We want to find what carries the most information
\vfill
For this, we are going to project the information in a new basis in which the first ``dimension'' will carry most information (in a sense we'll define later), the second dimension will carry a little less, etc.
\vfill
In order to do so, we need to learn how to change bases
\end{frame}

%%%%%%%%%%%%%%%%%%%%%
%%%%%%%%%%%%%%%%%%%%%
%%%%%%%%%%%%%%%%%%%%%
%%%%%%%%%%%%%%%%%%%%%
\SSsection{Change of basis}

%%%%%%%%%%%%%%%%%%%%%
%%%%%%%%%%%%%%%%%%%%%
\SSsubsection{Theory}


\begin{frame}
	In the following, 
	\[
	[\bx]_\B
	\]
	denotes the coordinates of $\bx$ in the basis $\B=\{\bu_1,\ldots,\bu_n\}$, i.e., $[\bx]_\B=(b_1,\ldots,b_n)^T$ such that $\bx = b_1\bu_1+\cdots+b_n\bu_n$
	\vfill
	The aim of a change of basis is to express vectors in another coordinate system (another basis) $\C=\{\bv_1,\ldots,\bv_n\}$
	\vfill
	We do so by finding a matrix allowing to move from one basis to another
\end{frame}

\begin{frame}{Change of basis}
\begin{definition}[Change of basis matrix]
Let $\B=\{\bu_1,\ldots,\bu_n\}$ and $\C=\{\bv_1,\ldots,\bv_n\}$ be bases of vector space $V$
\vskip0.5cm
The \textbf{change of basis matrix} $P_{\C\leftarrow\B}\in\M_n$ from basis $\B$ to $\C$ is
\[
P_{\C\leftarrow\B}
=\left[
[\bu_1]_\C \cdots [\bu_n]_\C
\right]
\]
i.e., its columns are the coordinate vectors $[\bu_1]_\C,\ldots,[\bu_n]_\C$ of vectors in $\B$ with respect to $\C$
\end{definition}
\end{frame}

\begin{frame}{Example -- $\B$ is the canonical basis of $\IR^3$}
Let $\B=\{\bi,\bj,\bk\}$ be the canonical basis of $\IR^3$ and $\C=\{\bv_1,\bv_2,\bv_3\}$, with
\[
\bv_1 = \begin{pmatrix} 1 \\ 0 \\ 0 \end{pmatrix}, \quad
\bv_2 = \begin{pmatrix} 1 \\ 1 \\ 0 \end{pmatrix}, \quad
\bv_3 = \begin{pmatrix} 1 \\ 1 \\ 1 \end{pmatrix}
\]
\vfill
To find $P_{\C\leftarrow\B} = [[\bi]_\C \;\; [\bj]_\C \;\; [\bk]_\C]$, express each canonical vector in terms of $\C$:
\[
\bi = \bv_1, \quad
\bj = \bv_2 - \bv_1, \quad
\bk = \bv_3 - \bv_2
\]
\[
\implies [\bi]_\C = \begin{pmatrix} 1 \\ 0 \\ 0 \end{pmatrix}, \quad
[\bj]_\C = \begin{pmatrix} -1 \\ 1 \\ 0 \end{pmatrix}, \quad
[\bk]_\C = \begin{pmatrix} 0 \\ -1 \\ 1 \end{pmatrix}
\]
\vfill
The change of basis matrix from $\B$ to $\C$ is
\[
P_{\C\leftarrow\B} = [[\bi]_\C \;\; [\bj]_\C \;\; [\bk]_\C] = \begin{pmatrix} 1 & -1 & 0 \\ 0 & 1 & -1 \\ 0 & 0 & 1 \end{pmatrix}
\]
\end{frame}


\begin{frame}{Example -- $\C$ is the canonical basis of $\IR^3$}
Now consider the reverse: let $\B=\{\bu_1,\bu_2,\bu_3\}$ with
\[
\bu_1 = \begin{pmatrix} 1 \\ 0 \\ 0 \end{pmatrix}, \quad
\bu_2 = \begin{pmatrix} 1 \\ 1 \\ 0 \end{pmatrix}, \quad
\bu_3 = \begin{pmatrix} 1 \\ 1 \\ 1 \end{pmatrix}
\]
and let $\C=\{\bi,\bj,\bk\}$ be the canonical basis of $\IR^3$
\vfill
Since $\C$ is the canonical basis, the coordinates in $\C$ are simply the vector itself:
\[
[\bu_1]_\C = \begin{pmatrix} 1 \\ 0 \\ 0 \end{pmatrix}, \quad
[\bu_2]_\C = \begin{pmatrix} 1 \\ 1 \\ 0 \end{pmatrix}, \quad
[\bu_3]_\C = \begin{pmatrix} 1 \\ 1 \\ 1 \end{pmatrix}
\]
\vfill
Thus, the change of basis matrix from $\B$ to the canonical basis $\C$ is
\[
P_{\C\leftarrow\B} = [[\bu_1]_\C \;\; [\bu_2]_\C \;\; [\bu_3]_\C] = \begin{pmatrix} 1 & 1 & 1 \\ 0 & 1 & 1 \\ 0 & 0 & 1 \end{pmatrix}
\]
\end{frame}

\begin{frame}
\begin{theorem}\label{th:change-of-basis-properties}
$\B=\{\bu_1,\ldots,\bu_n\}$ and $\C=\{\bv_1,\ldots,\bv_n\}$ bases of vector space $V$ and $P_{\C\leftarrow\B}$ a change of basis matrix from $\B$ to $\C$
\begin{enumerate}
\item $\forall\bx\in V$, $P_{\C\leftarrow\B}[\bx]_\B = [\bx]_\C$
\item $P_{\C\leftarrow\B}$ s.t. $\forall\bx\in V$, $P_{\C\leftarrow\B}[\bx]_\B = [\bx]_\C$ is \textbf{unique}
\item $P_{\C\leftarrow\B}$ invertible and $P_{\C\leftarrow\B}^{-1}=P_{\B\leftarrow\C}$
\end{enumerate}
\end{theorem}
\end{frame}

\begin{frame}{Example -- Linking the two directions}
Let $\B=\{\bi,\bj,\bk\}$ be the canonical basis and $\C=\{\bv_1,\bv_2,\bv_3\}$ with
\[
\bv_1 = \begin{pmatrix} 1 \\ 0 \\ 0 \end{pmatrix}, \quad
\bv_2 = \begin{pmatrix} 1 \\ 1 \\ 0 \end{pmatrix}, \quad
\bv_3 = \begin{pmatrix} 1 \\ 1 \\ 1 \end{pmatrix}
\]
\vfill
From our two previous examples:
\[
P_{\C\leftarrow\B} = \begin{pmatrix} 1 & -1 & 0 \\ 0 & 1 & -1 \\ 0 & 0 & 1 \end{pmatrix}
\quad\text{and}\quad
P_{\B\leftarrow\C} = \begin{pmatrix} 1 & 1 & 1 \\ 0 & 1 & 1 \\ 0 & 0 & 1 \end{pmatrix}
\]
\vfill
By item~3 in Theorem~\ref{th:change-of-basis-properties}, $P_{\C\leftarrow\B}^{-1} = P_{\B\leftarrow\C}$. Indeed:
\[
P_{\C\leftarrow\B} P_{\B\leftarrow\C}
= \begin{pmatrix} 1 & -1 & 0 \\ 0 & 1 & -1 \\ 0 & 0 & 1 \end{pmatrix}
\begin{pmatrix} 1 & 1 & 1 \\ 0 & 1 & 1 \\ 0 & 0 & 1 \end{pmatrix}
= \begin{pmatrix} 1 & 0 & 0 \\ 0 & 1 & 0 \\ 0 & 0 & 1 \end{pmatrix} = \II
\]
\vfill
Since $P_{\B\leftarrow\C} = [\bv_1 \;\; \bv_2 \;\; \bv_3]$ is immediate, we have $P_{\C\leftarrow\B} = [\bv_1 \;\; \bv_2 \;\; \bv_3]^{-1}$
\end{frame}


\begin{frame}{Row-reduction method for changing bases}
\begin{importanttheorem}
\label{th:change-basis-construction}
$\B=\{\bu_1,\ldots,\bu_n\}$ and $\C=\{\bv_1,\ldots,\bv_n\}$ bases of vector space $V$. Let $\E$ be any basis for $V$,
\[
B = [[\bu_1]_\E,\ldots,[\bu_n]_\E] 
\textrm{ and }
C = [[\bv_1]_\E,\ldots,[\bv_n]_\E] 
\]
and let $[C|B]$ be the augmented matrix constructed using $C$ and $B$. Then
\[
RREF\left([C|B]\right)
=[\II|P_{\C\leftarrow\B}]
\]
\end{importanttheorem}
\vfill
If working in $\IR^n$, this is quite useful with $\E$ the standard basis of $\IR^n$ (it does not matter if $\B=\E$)
\end{frame}

\begin{frame}
So the question now becomes
\begin{quote}
How do we find what new basis to look at our data in?
\end{quote}
\vfill
(Changing the basis does not change the data, just the view you have of it)
\vfill
(Think of what happens when you do a headstand.. your up becomes down, your right and left switch, but the world does not change, just your view of it)
\vfill
(Changes of bases are \emph{fundamental} operations in Science)
\end{frame}







%%%%%%%%%%%%%%%%%%%%%
%%%%%%%%%%%%%%%%%%%%%
\SSsubsection{A simple example}

\begin{frame}{Example: change of basis}
\begin{columns}[c]
\begin{column}{0.58\textwidth}
\textbf{Problem:} Find the change of basis matrix from basis $\B$ to basis $\C$ in $\IR^2$, where
\[
\B = \left\{\begin{pmatrix} 1 \\ 0 \end{pmatrix}, \begin{pmatrix} 0 \\ 1 \end{pmatrix}\right\}
\quad\text{and}\quad
\C = \left\{\begin{pmatrix} 1 \\ 1 \end{pmatrix}, \begin{pmatrix} 1 \\ -1 \end{pmatrix}\right\}
\]
\vskip0.5cm
Then use this matrix to find the coordinates of $\bx = \begin{pmatrix} 3 \\ 2 \end{pmatrix}$ in basis $\C$
\end{column}
\begin{column}{0.42\textwidth}
\centering
\begin{tikzpicture}[scale=0.85, >=Latex]
    % Grid
    \draw[step=1cm, gray!20, very thin] (-1.2,-1.8) grid (3.8,2.8);
    % Axes
    \draw[->, thick, gray!60] (-1.2,0) -- (3.8,0) node[right, font=\footnotesize] {$x$};
    \draw[->, thick, gray!60] (0,-1.8) -- (0,2.8) node[above, font=\footnotesize] {$y$};
    
    % Basis B vectors (standard)
    \draw[->, very thick, black] (0,0) -- (1,0) node[below=2pt, font=\small] {$\bi$};
    \draw[->, very thick, black] (0,0) -- (0,1) node[left=2pt, font=\small] {$\bj$};
    
    % New axes directions (dashed)
    \draw[dashed, blue!50, thin] (-1.5,-1.5) -- (2.8,2.8) node[above right, font=\footnotesize, text=blue!80!black] {$\bv_1$-axis};
    \draw[dashed, blue!50, thin] (-1.5,1.5) -- (2.3,-2.3) node[below right, font=\footnotesize, text=blue!80!black] {$\bv_2$-axis};
    
    % Basis C vectors
    \draw[->, very thick, blue!80!black] (0,0) -- (1,1) node[above left, font=\small] {$\bv_1$};
    \draw[->, very thick, blue!80!black] (0,0) -- (1,-1) node[below left, font=\small] {$\bv_2$};
    
    % Vector x
    \draw[->, very thick, red!80!black] (0,0) -- (3,2) node[above, font=\small] {$\bx$};
    \filldraw[red!80!black] (3,2) circle (1.5pt);
\end{tikzpicture}
\end{column}
\end{columns}
\end{frame}

\begin{frame}{Step 1 -- Set up the matrices}
\begin{itemize}
\item $\B=\{\bi,\bj\}$ is the standard basis of $\IR^2$, so using $\E=\B$, we have 
\[
[\bi]_\E = \begin{pmatrix} 1 \\ 0 \end{pmatrix}, \quad [\bj]_\E = \begin{pmatrix} 0 \\ 1 \end{pmatrix}
\]
\item For $\C$: 
\[
[\bv_1]_\E = \begin{pmatrix} 1 \\ 1 \end{pmatrix}, \quad [\bv_2]_\E = \begin{pmatrix} 1 \\ -1 \end{pmatrix}
\]
\end{itemize}
\end{frame}

\begin{frame}{Step 2 -- Row reduce}
Using Theorem~\ref{th:change-basis-construction}, we form the augmented matrix $[C|B]$:
$$[C|B] = \left[\begin{array}{cc|cc} 1 & 1 & 1 & 0 \\ 1 & -1 & 0 & 1 \end{array}\right]$$

\vfill
\textbf{Row reduce to RREF:}
\begin{align*}
\left[\begin{array}{cc|cc} 1 & 1 & 1 & 0 \\ 1 & -1 & 0 & 1 \end{array}\right] &\xrightarrow{R_2 \leftarrow R_2 - R_1} \left[\begin{array}{cc|cc} 1 & 1 & 1 & 0 \\ 0 & -2 & -1 & 1 \end{array}\right] \\
&\xrightarrow{R_2 \leftarrow -\frac{1}{2}R_2} \left[\begin{array}{cc|cc} 1 & 1 & 1 & 0 \\ 0 & 1 & \frac{1}{2} & -\frac{1}{2} \end{array}\right] \\
&\xrightarrow{R_1 \leftarrow R_1 - R_2} \left[\begin{array}{cc|cc} 1 & 0 & \frac{1}{2} & \frac{1}{2} \\ 0 & 1 & \frac{1}{2} & -\frac{1}{2} \end{array}\right]
\end{align*}
\end{frame}

\begin{frame}{Step 3 -- Extract the change of basis matrix}
From the RREF form $[\II|P_{\C\leftarrow\B}]$, we get:
$$P_{\C\leftarrow\B} = \begin{pmatrix} \frac{1}{2} & \frac{1}{2} \\ \frac{1}{2} & -\frac{1}{2} \end{pmatrix}$$

\vfill
\textbf{Verification:} Let's check that this matrix works correctly.
\begin{itemize}
\item $[\bi]_\C = P_{\C\leftarrow\B}[\bi]_\B = \begin{pmatrix} \frac{1}{2} & \frac{1}{2} \\ \frac{1}{2} & -\frac{1}{2} \end{pmatrix}\begin{pmatrix} 1 \\ 0 \end{pmatrix} = \begin{pmatrix} \frac{1}{2} \\ \frac{1}{2} \end{pmatrix}$
\item Check: $\frac{1}{2}\begin{pmatrix} 1 \\ 1 \end{pmatrix} + \frac{1}{2}\begin{pmatrix} 1 \\ -1 \end{pmatrix} = \begin{pmatrix} 1 \\ 0 \end{pmatrix}$
\end{itemize}
\end{frame}

\begin{frame}{Step 4 -- Coordinates of $\bx$ in basis $\C$}
Now we find $[\bx]_\C$ for $\bx = \begin{pmatrix} 3 \\ 2 \end{pmatrix}$:

$$[\bx]_\C = P_{\C\leftarrow\B}[\bx]_\B = \begin{pmatrix} \frac{1}{2} & \frac{1}{2} \\ \frac{1}{2} & -\frac{1}{2} \end{pmatrix}\begin{pmatrix} 3 \\ 2 \end{pmatrix} = \begin{pmatrix} \frac{5}{2} \\ \frac{1}{2} \end{pmatrix}$$

\vfill
\textbf{Verification:} Check that this gives us back the original vector:
$$\frac{5}{2}\begin{pmatrix} 1 \\ 1 \end{pmatrix} + \frac{1}{2}\begin{pmatrix} 1 \\ -1 \end{pmatrix} = \begin{pmatrix} \frac{5}{2} + \frac{1}{2} \\ \frac{5}{2} - \frac{1}{2} \end{pmatrix} = \begin{pmatrix} 3 \\ 2 \end{pmatrix}$$

\vfill
\textbf{Answer:} The coordinates of $\bx = \begin{pmatrix} 3 \\ 2 \end{pmatrix}$ in basis $\C$ are $[\bx]_\C = \begin{pmatrix} \frac{5}{2} \\ \frac{1}{2} \end{pmatrix}$
\end{frame}

\begin{frame}{Alternative method -- Direct calculation}
We could also solve this directly by setting up the system:
$$\bx = c_1\bv_1 + c_2\bv_2$$
$$\begin{pmatrix} 3 \\ 2 \end{pmatrix} = c_1\begin{pmatrix} 1 \\ 1 \end{pmatrix} + c_2\begin{pmatrix} 1 \\ -1 \end{pmatrix}$$

This gives us the system:
\begin{align*}
c_1 + c_2 &= 3 \\
c_1 - c_2 &= 2
\end{align*}

Solving: $c_1 = \frac{5}{2}$, $c_2 = \frac{1}{2}$

\vfill
This confirms our result: $[\bx]_\C = \begin{pmatrix} \frac{5}{2} \\ \frac{1}{2} \end{pmatrix}$
\end{frame}

\begin{frame}[fragile]{Computation in R -- Finding $P_{\C\leftarrow\B}$}
Set up $C = [\bv_1, \bv_2]$, $B = [\bi, \bj]$, form the augmented matrix $[C|B]$ and row reduce:
\vfill
\begin{knitrout}
\definecolor{shadecolor}{rgb}{0.969, 0.969, 0.969}\color{fgcolor}\begin{kframe}
\begin{alltt}
\hldef{C} \hlkwb{=} \hlkwd{matrix}\hldef{(}\hlkwd{c}\hldef{(}\hlnum{1}\hldef{,} \hlnum{1}\hldef{,} \hlnum{1}\hldef{,} \hlopt{-}\hlnum{1}\hldef{),} \hlkwc{nrow} \hldef{=} \hlnum{2}\hldef{)}
\hldef{B} \hlkwb{=} \hlkwd{diag}\hldef{(}\hlnum{2}\hldef{)}
\hldef{CB} \hlkwb{=} \hlkwd{cbind}\hldef{(C, B)}
\hldef{(R} \hlkwb{=} \hldef{pracma}\hlopt{::}\hlkwd{rref}\hldef{(CB))}
\end{alltt}
\begin{verbatim}
##      [,1] [,2] [,3] [,4]
## [1,]    1    0  0.5  0.5
## [2,]    0    1  0.5 -0.5
\end{verbatim}
\begin{alltt}
\hldef{(P} \hlkwb{=} \hldef{R[,} \hlnum{3}\hlopt{:}\hlnum{4}\hldef{])}
\end{alltt}
\begin{verbatim}
##      [,1] [,2]
## [1,]  0.5  0.5
## [2,]  0.5 -0.5
\end{verbatim}
\end{kframe}
\end{knitrout}
\end{frame}

\begin{frame}[fragile]{Computation in R -- Coordinates of $\bx$}
Now compute $[\bx]_\C = P_{\C\leftarrow\B} [\bx]_\B$ and verify that $C [\bx]_\C = \bx$:
\vfill
\begin{knitrout}
\definecolor{shadecolor}{rgb}{0.969, 0.969, 0.969}\color{fgcolor}\begin{kframe}
\begin{alltt}
\hldef{x} \hlkwb{=} \hlkwd{c}\hldef{(}\hlnum{3}\hldef{,} \hlnum{2}\hldef{)}
\hldef{(x_C} \hlkwb{=} \hldef{P} \hlopt{%*%} \hldef{x)}
\end{alltt}
\begin{verbatim}
##      [,1]
## [1,]  2.5
## [2,]  0.5
\end{verbatim}
\begin{alltt}
\hlcom{# Verification: reconstruct x in the canonical basis}
\hldef{C} \hlopt{%*%} \hldef{x_C}
\end{alltt}
\begin{verbatim}
##      [,1]
## [1,]    3
## [2,]    2
\end{verbatim}
\end{kframe}
\end{knitrout}
\end{frame}

%%%%%%%%%%%%%%%%%%%%%
%%%%%%%%%%%%%%%%%%%%%
\SSsubsection{Basis for a trajectory in \texorpdfstring{$\IR^2$}{R\textasciicircum 2}}

\begin{frame}{Example -- A trajectory in $\IR^2$}
\begin{columns}[c]
\begin{column}{0.54\textwidth}
Let $\B=\{\bi,\bj\}$ be the canonical basis of $\IR^2$.
\vskip0.2cm
Consider a trajectory along the unit circle:
\[
\bx(t) = \begin{pmatrix} \cos t \\ \sin t \end{pmatrix}, \quad t \in \IR
\]
\vskip0.2cm
This is a function $\bx:\IR\to\IR^2$ that we call a \textbf{parametric equation}:
\begin{itemize}
    \item $t$ is the parameter (often interpreted as time)
    \item For each $t$, $\bx(t)$ is a point in $\IR^2$; as $t$ varies, $\bx(t)$ traces out the trajectory
\end{itemize}
\end{column}
\begin{column}{0.46\textwidth}
\centering
\begin{tikzpicture}[scale=0.95, >=Latex]
    % Axes
    \draw[->, gray!40] (-2.2,0) -- (2.4,0) node[right, font=\tiny] {$x$};
    \draw[->, gray!40] (0,-2.2) -- (0,2.4) node[above, font=\tiny] {$y$};
    
    % Circle trajectory
    \draw[thick, gray!50] (0,0) circle (1.6);
    
    % Standard basis at origin
    \draw[->, very thick, black] (0,0) -- (0.8,0) node[below=1pt, font=\scriptsize] {$\bi$};
    \draw[->, very thick, black] (0,0) -- (0,0.8) node[left=1pt, font=\scriptsize] {$\bj$};
    
    % Direction arrow on trajectory
    \draw[->, thick, teal] (15:1.6) arc (15:35:1.6);
    \draw[->, thick, teal] (105:1.6) arc (105:125:1.6);

    % Point 1 at t = 45 deg
    \coordinate (P1) at (45:1.6);
    \draw[->, thick, purple!70!black] (0,0) -- (P1) node[midway, below right, font=\scriptsize] {$\bx(t_1)$};
    \filldraw[purple!80!black] (P1) circle (1.8pt);
    
    % Point 2 at t = 135 deg
    \coordinate (P2) at (135:1.6);
    \draw[->, thick, purple!70!black] (0,0) -- (P2) node[midway, below left, font=\scriptsize] {$\bx(t_2)$};
    \filldraw[purple!80!black] (P2) circle (1.8pt);
\end{tikzpicture}
\end{column}
\end{columns}
\end{frame}

\begin{frame}{A moving basis along the trajectory}
\begin{columns}[c]
\begin{column}{0.54\textwidth}
At each point $\bx(t)$ on the trajectory, define a moving basis $\C(t) = \{\bv_1(t), \bv_2(t)\}$:
\begin{itemize}
    \item $\bv_1(t) = \bx'\!(t) = \begin{pmatrix} -\sin t \\ \cos t \end{pmatrix}$ (unit tangent)
    \item $\bv_2(t) = \begin{pmatrix} \cos t \\ \sin t \end{pmatrix}$ (outward unit normal)
\end{itemize}
\vskip0.2cm
Since $\bv_1(t) \bullet \bv_2(t) = 0$ and $\|\bv_1\|=\|\bv_2\|=1$, $\C(t)$ is an orthonormal basis for all $t$
\vskip0.2cm
As $t$ changes, the basis $\{\bv_1(t), \bv_2(t)\}$ moves and rotates along with $\bx(t)$
\end{column}
\begin{column}{0.46\textwidth}
\centering
\begin{tikzpicture}[scale=0.95, >=Latex]
    % Axes
    \draw[->, gray!40] (-2.2,0) -- (2.4,0) node[right, font=\tiny] {$x$};
    \draw[->, gray!40] (0,-2.2) -- (0,2.4) node[above, font=\tiny] {$y$};
    
    % Circle trajectory
    \draw[thick, gray!50] (0,0) circle (1.6);
    
    % Standard basis at origin
    \draw[->, very thick, black] (0,0) -- (0.8,0) node[below=1pt, font=\scriptsize] {$\bi$};
    \draw[->, very thick, black] (0,0) -- (0,0.8) node[left=1pt, font=\scriptsize] {$\bj$};
    
    % Direction arrow on trajectory
    \draw[->, thick, teal] (15:1.6) arc (15:35:1.6);
    \draw[->, thick, teal] (105:1.6) arc (105:125:1.6);

    % Point 1 at t = 45 deg
    \coordinate (P1) at (45:1.6);
    \draw[->, thick, purple!70!black] (0,0) -- (P1) node[midway, below right, font=\scriptsize] {$\bx(t_1)$};
    \filldraw[purple!80!black] (P1) circle (1.8pt);
    \draw[->, very thick, blue!80!black] (P1) -- ++(135:1.1) node[above left, font=\scriptsize] {$\bv_1(t_1)$};
    \draw[->, very thick, red!80!black] (P1) -- ++(45:1.1) node[above right, font=\scriptsize] {$\bv_2(t_1)$};
    
    % Point 2 at t = 135 deg
    \coordinate (P2) at (135:1.6);
    \draw[->, thick, purple!70!black] (0,0) -- (P2) node[midway, below left, font=\scriptsize] {$\bx(t_2)$};
    \filldraw[purple!80!black] (P2) circle (1.8pt);
    \draw[->, very thick, blue!80!black] (P2) -- ++(225:1.1) node[below left, font=\scriptsize] {$\bv_1(t_2)$};
    \draw[->, very thick, red!80!black] (P2) -- ++(135:1.1) node[above left, font=\scriptsize] {$\bv_2(t_2)$};
\end{tikzpicture}
\end{column}
\end{columns}
\end{frame}

\begin{frame}{Coordinates in the moving basis}
The change of basis matrix from $\C(t)$ to the canonical basis $\B$ is
\[
P_{\B\leftarrow\C(t)} = [\bv_1(t) \;\; \bv_2(t)] = \begin{pmatrix} -\sin t & \cos t \\ \cos t & \sin t \end{pmatrix}
\]
\vfill
Since $P_{\B\leftarrow\C(t)}$ is symmetric and orthogonal, the change of basis matrix to $\C(t)$ is:
\[
P_{\C(t)\leftarrow\B} = P_{\B\leftarrow\C(t)}^{-1} = P_{\B\leftarrow\C(t)}^T = \begin{pmatrix} -\sin t & \cos t \\ \cos t & \sin t \end{pmatrix}
\]
\vfill
In the moving basis $\C(t)$, the coordinates of position and velocity are:
\[
[\bx(t)]_{\C(t)} = P_{\C(t)\leftarrow\B} \begin{pmatrix} \cos t \\ \sin t \end{pmatrix} = \begin{pmatrix} 0 \\ 1 \end{pmatrix}, \qquad
[\bx'\!(t)]_{\C(t)} = P_{\C(t)\leftarrow\B} \begin{pmatrix} -\sin t \\ \cos t \end{pmatrix} = \begin{pmatrix} 1 \\ 0 \end{pmatrix}
\]
\vfill
In the local frame, the particle is always at distance $1$ along the normal axis with velocity $1$ along the tangent
\vfill
\footnotesize\emph{Choosing the inward normal $-\bv_2(t)$ simply gives $[\bx(t)]_{\C(t)} = \begin{pmatrix} 0 \\ -1 \end{pmatrix}$}
\end{frame}


%%%%%%%%%%%%%%%%%%%%%
%%%%%%%%%%%%%%%%%%%%%
\SSsection{A crash course on probability}


\begin{frame}{Why probability?}
We said earlier that we would look for a basis in which the first dimension carries most information
\vfill
But how do we define \emph{information}?
\vfill
We use concepts from probability and statistics to do so
\vfill
A good measure of information is \emph{variance} (how much data varies around the mean)
\end{frame} 

\begin{frame}{Brief ``review'' of some probability concepts}
Proper definition of \emph{probability} requires to use \emph{measure theory}.. will not get into details here
\vfill
A \textbf{random variable} $X$ is a \emph{measurable} function $X:\Omega\to E$, where $\Omega$ is a set of outcomes (\emph{sample space}) and $E$ is a measurable space
\vfill
$\IP(X\in S\subseteq E) = \IP(\omega\in\Omega|X(\omega)\in S)$
\vfill
\textbf{Distribution function} of a r.v., $F(x)=\IP(X\leq x)$, describes the distribution of a r.v.
\vfill
R.v. can be discrete or continuous or .. other things. 
\end{frame}

\begin{frame}{Variance}
\begin{definition}[Variance]
Let $X$ be a random variable. The \textbf{variance} of $X$ is given by
\[
\Var X = E\left[\left(X-E(X)\right)^2\right]
\]
where $E$ is the expected value
\end{definition}
\end{frame}

\begin{frame}{Covariance}
\begin{definition}[Covariance]
Let $X,Y$ be jointly distributed random variables. The \textbf{covariance} of $X$ and $Y$ is given by
\[
\cov (X,Y) = E\left[\left(X-E(X)\right)\left(Y-E(Y)\right)\right]
\]
\end{definition}
\vfill
Note that $\cov(X,X)=E\left[\left(X-E(X)\right)^2\right] = \Var X$
\end{frame}

\begin{frame}{In practice: ``true law'' versus ``observation''}
In statistics: we reason on the \emph{true law} of distributions, but we usually have only access to a sample
\vfill
We then use \textbf{estimators} to .. estimate the value of a parameter, e.g., the mean, variance and covariance
\vfill
\end{frame}
    
\begin{frame}
\begin{definition}[Unbiased estimators of the mean and variance]
Let $x_1,\ldots,x_n$ be data points (the \emph{sample}) and 
\[
\bar x = \frac 1n \sum_{i=1}^n x_i
\]
be the \textbf{mean} of the data. An unbiased estimator of the variance of the sample is
\[
\sigma^2 = \frac{1}{n-1}\sum_{i=1}^n (x_i-\bar x)^2
\]
\end{definition}
\vfill
\textbf{Unbiased}: the expected value of the estimator equals the true parameter value ($E(\bar x) = E(X)$ and $E(\sigma^2) = \Var X$): on average, it does not systematically overestimate or underestimate the parameter
\end{frame}

\begin{frame}[fragile]{Variance of hockey players' height and weight}
We compute the sample variance in R using the function \code{var}:
\vfill
\begin{knitrout}
\definecolor{shadecolor}{rgb}{0.969, 0.969, 0.969}\color{fgcolor}\begin{kframe}
\begin{alltt}
\hlkwd{var}\hldef{(data}\hlopt{$}\hldef{height)}
\end{alltt}
\begin{verbatim}
## [1] 29.66176
\end{verbatim}
\begin{alltt}
\hlkwd{var}\hldef{(data}\hlopt{$}\hldef{weight)}
\end{alltt}
\begin{verbatim}
## [1] 47.81112
\end{verbatim}
\end{kframe}
\end{knitrout}
\vfill
The variance is about $29.7\text{ cm}^2$ for height and $47.8\text{ kg}^2$ for weight
\end{frame}

\begin{frame}[fragile]{Normalising does not change the variance}
Centring (normalising) shifts data by a constant, leaving spread unchanged:
\[
\Var(X - c) = \Var X
\]
\vfill
Indeed, computing the variance using the centred values:
\vfill
\begin{knitrout}
\definecolor{shadecolor}{rgb}{0.969, 0.969, 0.969}\color{fgcolor}\begin{kframe}
\begin{alltt}
\hlkwd{var}\hldef{(data}\hlopt{$}\hldef{height.c)}
\end{alltt}
\begin{verbatim}
## [1] 29.66176
\end{verbatim}
\begin{alltt}
\hlkwd{var}\hldef{(data}\hlopt{$}\hldef{weight.c)}
\end{alltt}
\begin{verbatim}
## [1] 47.81112
\end{verbatim}
\end{kframe}
\end{knitrout}
\vfill
The variance is identical to that of the uncentred data
\end{frame}

\begin{frame}
\begin{definition}[Unbiased estimator of the covariance - 2D case]
Let $(x_1,y_1),\ldots,(x_n,y_n)$ be data points,
\[
\bar x = \frac 1n \sum_{i=1}^n x_i
\textrm{ and }
\bar y = \frac 1n \sum_{i=1}^n y_i
\]
be the means of the data. An unbiased estimator of the covariance of the sample is
\[
\cov(x,y) = \frac{1}{n-1}\sum_{i=1}^n (x_i-\bar x)(y_i-\bar y)
\]
\end{definition}
\end{frame}

\begin{frame}[fragile]{Computation of covariance in R -- ``By hand'' and with a function}
We compute the means and sum directly from the definition:
\vfill
\begin{knitrout}
\definecolor{shadecolor}{rgb}{0.969, 0.969, 0.969}\color{fgcolor}\begin{kframe}
\begin{alltt}
\hldef{n} \hlkwb{=} \hlkwd{length}\hldef{(data}\hlopt{$}\hldef{height)}
\hldef{x_bar} \hlkwb{=} \hlkwd{sum}\hldef{(data}\hlopt{$}\hldef{height)} \hlopt{/} \hldef{n}
\hldef{y_bar} \hlkwb{=} \hlkwd{sum}\hldef{(data}\hlopt{$}\hldef{weight)} \hlopt{/} \hldef{n}
\hlkwd{sum}\hldef{((data}\hlopt{$}\hldef{height} \hlopt{-} \hldef{x_bar)} \hlopt{*} \hldef{(data}\hlopt{$}\hldef{weight} \hlopt{-} \hldef{y_bar))} \hlopt{/} \hldef{(n} \hlopt{-} \hlnum{1}\hldef{)}
\end{alltt}
\begin{verbatim}
## [1] 26.63506
\end{verbatim}
\begin{alltt}
\hlkwd{cov}\hldef{(data}\hlopt{$}\hldef{height, data}\hlopt{$}\hldef{weight)}
\end{alltt}
\begin{verbatim}
## [1] 26.63506
\end{verbatim}
\end{kframe}
\end{knitrout}
\vfill
The covariance is about $26.6\text{ cm}\cdot\text{kg}$, matching R's built-in \code{cov}
\end{frame}

\begin{frame}{What does covariance do?}
Variance explains how data disperses around the mean, in a 1-D case
\vfill
Covariance measures the relationship between two dimensions. E.g., height and weight
\vfill
More than the exact value, the sign is important:
\begin{itemize}
    \item $\cov(X,Y)>0$: both dimensions change in the same ``direction''; e.g., larger height usually means higher weight
    \item $\cov(X,Y)<0$: both dimensions change in reverse directions; e.g., time spent on social media and performance in this class
    \item $\cov(X,Y)=0$: the dimensions are independent from one another; e.g., sex/gender and ``intelligence''
\end{itemize}
\end{frame}

\begin{frame}{The covariance matrix}
\begin{definition}[Covariance matrix]
Suppose $p$ random variables $X_1,\ldots,X_p$. Then the covariance matrix is the symmetric matrix
\[
\begin{pmatrix}
\Var X_1 & \cov(X_1,X_2) & \cdots & \cov(X_1,X_p) \\
\cov(X_1,X_2) & \Var X_2 & \cdots & \cov(X_2,X_p) \\
\vdots & \vdots & & \vdots \\
\cov(X_1,X_p) & \cov(X_2,X_p) & \cdots & \Var X_p 
\end{pmatrix}
\]
\end{definition}
\vfill
The actual form of the matrix is below but the form above stems from the properties of covariance
\[
\begin{pmatrix}
\cov(X_1,X_1) & \cov(X_1,X_2) & \cdots & \cov(X_1,X_p) \\
\cov(X_2,X_1) & \cov(X_2,X_2) & \cdots & \cov(X_2,X_p) \\
\vdots & \vdots & & \vdots \\
\cov(X_p,X_1) & \cov(X_p,X_2) & \cdots & \cov(X_p,X_p) 
\end{pmatrix}
\]
\end{frame}

\begin{frame}{Sample covariance matrix}
Let $X\in\M_{np}$ be the data matrix ($n$ observations in rows, $p$ features or variables in columns) and $\widetilde{X}\in\M_{np}$ be the \textbf{centred data matrix}. Then the (sample) covariance matrix $S_X\in\M_p$ takes the form
\[
S_X = \frac{1}{n-1}\widetilde{X}^T\widetilde{X}
\]
where the data is centred!
\vfill
Sometimes you will see $S_X=1/(n-1)\widetilde{X}\widetilde{X}^T$. This is for matrices with observations in columns and features (variables) in rows. Just remember that you want the covariance matrix to have size the number of features ($p$), not observations ($n$), this will give you the order in which to take the product
\end{frame}

\begin{frame}{Why is $S_X = \frac{1}{n-1}\widetilde{X}^T\widetilde{X}$? -- Data matrix setup}
Let us prove why the sample covariance matrix takes this matrix form
\vfill
Suppose we have $n$ observations of $p$ features (variables) $X_1,\ldots,X_p$
\vfill
Organise the centred data into the \defword{centred data matrix} $\widetilde{X}\in\M_{np}$ where each row is an observation and each column is a centred feature:
\[
\widetilde{X} = \begin{pmatrix}
\tilde{x}_{11} & \tilde{x}_{12} & \cdots & \tilde{x}_{1p} \\
\tilde{x}_{21} & \tilde{x}_{22} & \cdots & \tilde{x}_{2p} \\
\vdots & \vdots & \ddots & \vdots \\
\tilde{x}_{n1} & \tilde{x}_{n2} & \cdots & \tilde{x}_{np}
\end{pmatrix}
= \begin{pmatrix} \tilde{\bc}_1 & \tilde{\bc}_2 & \cdots & \tilde{\bc}_p \end{pmatrix}
\]
where $\tilde{x}_{ij} = x_{ij} - \bar{x}_j$ and column $\tilde{\bc}_j = (\tilde{x}_{1j}, \tilde{x}_{2j}, \ldots, \tilde{x}_{nj})^T \in \IR^n$ contains all $n$ centred measurements of feature $X_j$
\vfill
Because the data is \defword{centred}, $\frac{1}{n}\sum_{i=1}^n \tilde{x}_{ij} = 0$ for all $j=1,\ldots,p$
\end{frame}

\begin{frame}{Why is $S_X = \frac{1}{n-1}\widetilde{X}^T\widetilde{X}$? -- Sample covariances}
Recall the definition of the sample covariance between features $X_j$ and $X_k$:
\[
\cov(X_j, X_k) = \frac{1}{n-1} \sum_{i=1}^n (x_{ij} - \bar x_j)(x_{ik} - \bar x_k)
\]
\vfill
Writing $\tilde{x}_{ij} = x_{ij} - \bar x_j$ and $\tilde{x}_{ik} = x_{ik} - \bar x_k$, this simplifies to:
\[
\cov(X_j, X_k) = \frac{1}{n-1} \sum_{i=1}^n \tilde{x}_{ij} \tilde{x}_{ik}
\]
\vfill
In particular, when $j = k$, we obtain the sample variance of feature $X_j$:
\[
\Var X_j = \cov(X_j, X_j) = \frac{1}{n-1} \sum_{i=1}^n \tilde{x}_{ij}^2
\]
\vfill
Notice: $\sum_{i=1}^n \tilde{x}_{ij} \tilde{x}_{ik} = \tilde{\bc}_j^T \tilde{\bc}_k$, the dot product of column $\tilde{\bc}_j$ and column $\tilde{\bc}_k$!
\end{frame}

\begin{frame}{Why is $S_X = \frac{1}{n-1}\widetilde{X}^T\widetilde{X}$? -- Matrix product $\widetilde{X}^T\widetilde{X}$}
Now consider the matrix product $\widetilde{X}^T\widetilde{X}$:
\vfill
$\widetilde{X}^T \in \M_{pn}$ has the columns of $\widetilde{X}$ as its rows:
\[
\widetilde{X}^T = \begin{pmatrix} \tilde{\bc}_1^T \\ \tilde{\bc}_2^T \\ \vdots \\ \tilde{\bc}_p^T \end{pmatrix}
= \begin{pmatrix}
\tilde{x}_{11} & \tilde{x}_{21} & \cdots & \tilde{x}_{n1} \\
\tilde{x}_{12} & \tilde{x}_{22} & \cdots & \tilde{x}_{n2} \\
\vdots & \vdots & \ddots & \vdots \\
\tilde{x}_{1p} & \tilde{x}_{2p} & \cdots & \tilde{x}_{np}
\end{pmatrix}
\]
\vfill
Since $\widetilde{X}^T\in\M_{pn}$ and $\widetilde{X}\in\M_{np}$, the product $\widetilde{X}^T\widetilde{X}$ is a $p\times p$ matrix
\vfill
Its $(j,k)$-entry is the product of row $j$ of $\widetilde{X}^T$ and column $k$ of $\widetilde{X}$:
\[
(\widetilde{X}^T\widetilde{X})_{jk} = \tilde{\bc}_j^T\tilde{\bc}_k = \sum_{i=1}^n \tilde{x}_{ij}\tilde{x}_{ik}
\]
\end{frame}

\begin{frame}{Why is $S_X = \frac{1}{n-1}\widetilde{X}^T\widetilde{X}$? -- Conclusion}
Dividing by $n-1$, each entry of $\frac{1}{n-1}\widetilde{X}^T\widetilde{X}$ is:
\[
\left(\frac{1}{n-1}\widetilde{X}^T\widetilde{X}\right)_{jk} = \frac{1}{n-1}\sum_{i=1}^n \tilde{x}_{ij}\tilde{x}_{ik} = 
\begin{cases}
\Var X_j & \text{if } j = k \\
\cov(X_j, X_k) & \text{if } j \neq k
\end{cases}
\]
\vfill
Thus, writing out the matrix:
\[
\frac{1}{n-1}\widetilde{X}^T\widetilde{X} = \begin{pmatrix}
\Var X_1 & \cov(X_1, X_2) & \cdots & \cov(X_1, X_p) \\
\cov(X_2, X_1) & \Var X_2 & \cdots & \cov(X_2, X_p) \\
\vdots & \vdots & \ddots & \vdots \\
\cov(X_p, X_1) & \cov(X_p, X_2) & \cdots & \Var X_p
\end{pmatrix} = S_X
\]
\vfill
This proves that $S_X = \frac{1}{n-1}\widetilde{X}^T\widetilde{X}$ for the centred data matrix $\widetilde{X}$
\end{frame}

\begin{frame}{What if observations are in columns?}
Suppose the centred data is stored with observations in columns and features (variables) in rows:
\[
\widetilde{Y} = \widetilde{X}^T \in \M_{pn}
\]
where row $j$ of $\widetilde{Y}$ is $\tilde{\bc}_j^T = (\tilde{x}_{1j}, \ldots, \tilde{x}_{nj}) \in \IR^{1\times n}$ (feature $X_j$)
\vfill
Then $\widetilde{Y}\widetilde{Y}^T \in \M_p$, and its $(j,k)$-entry is:
\[
(\widetilde{Y}\widetilde{Y}^T)_{jk} = \tilde{\bc}_j^T\tilde{\bc}_k = \sum_{i=1}^n \tilde{x}_{ij}\tilde{x}_{ik}
\]
\vfill
Dividing by $n-1$:
\[
S_Y = \frac{1}{n-1}\widetilde{Y}\widetilde{Y}^T = S_X
\]
\vfill
The resulting sample covariance matrix is identical
\end{frame}

\begin{frame}{What if the data is not centred beforehand?}
Suppose $X\in\M_{np}$ is the \defword{uncentred} data matrix
\vfill
Let $\mathbf{1} = (1, 1, \ldots, 1)^T \in \IR^n$. The vector of column means is $\bar\bx^T = \frac{1}{n}\mathbf{1}^T X \in \IR^{1\times p}$
\vfill
Subtracting the means from each row of $X$ gives the \defword{centred data matrix} $\widetilde{X}\in\M_{np}$:
\[
\widetilde{X} = X - \mathbf{1}\bar\bx^T = X - \frac{1}{n}\mathbf{1}\mathbf{1}^T X = \left(\II_n - \frac{1}{n}\mathbf{1}\mathbf{1}^T\right) X = C_n X
\]
where $C_n = \II_n - \frac{1}{n}\mathbf{1}\mathbf{1}^T\in\M_n$ is the \defword{centring matrix}
\vfill
$C_n$ is symmetric ($C_n^T=C_n$) and idempotent ($C_n^2=C_n$), so:
\[
S_X = \frac{1}{n-1}\widetilde{X}^T\widetilde{X} = \frac{1}{n-1}X^T C_n^T C_n X = \frac{1}{n-1}X^T C_n X
\]
\end{frame}


\begin{frame}[fragile]{Covariance matrix of hockey players in R}
When passed a data frame or matrix, \code{cov} computes the covariance matrix:
\vfill
\begin{knitrout}
\definecolor{shadecolor}{rgb}{0.969, 0.969, 0.969}\color{fgcolor}\begin{kframe}
\begin{alltt}
\hlkwd{cov}\hldef{(data[,} \hlkwd{c}\hldef{(}\hlsng{"height"}\hldef{,} \hlsng{"weight"}\hldef{)])}
\end{alltt}
\begin{verbatim}
##          height   weight
## height 29.66176 26.63506
## weight 26.63506 47.81112
\end{verbatim}
\end{kframe}
\end{knitrout}
\vfill
The diagonal entries are the variances ($\Var(\text{height})\approx 29.7$, $\Var(\text{weight})\approx 47.8$) and the off-diagonal entries are $\cov(\text{height}, \text{weight})\approx 26.6$
\end{frame}


%%%%%%%%%%%%%%%%%%%%%
%%%%%%%%%%%%%%%%%%%%%
\SSsection{Back to PCA}


\begin{frame}{Setting things up}
I will use notation (mostly) as in Joliffe's \emph{Principal Component Analysis} (PDF of older version available for free from UofM Libraries)
\vfill
$\bx=(x_1,\ldots,x_p)$ vector of $p$ random variables
\end{frame}


\begin{frame} 
We seek a linear function $\bm{\alpha}_1^T\bx$ with maximum variance, where $\bm{\alpha}_1=(\alpha_{11},\ldots,\alpha_{1p})$, i.e.,
\[
\bm{\alpha}_1^T\bx = \sum_{j=1}^p\alpha_{1j}x_j
\]
\vfill
Then we seek a linear function $\bm{\alpha}_2^T\bx$ with maximum variance, uncorrelated to $\bm{\alpha}_1^T\bx$
\vfill
And we continue...
\vfill
At $k$th stage, we find a linear function $\bm{\alpha}_k^T\bx$ with maximum variance, uncorrelated to $\bm{\alpha}_1^T\bx,\ldots,\bm{\alpha}_{k-1}^T\bx$
\vfill
$\bm{\alpha}_i^T\bx$ is the $i$th \textbf{principal component} (PC)
\end{frame}

\begin{frame}{Case of known covariance matrix}
Suppose we know $\Sigma$, covariance matrix of $\bx$ (i.e., typically: we know $\bx$)
\vfill
Then the $k$th PC is 
\[
z_k=\bm{\alpha}_k^T\bx
\]
where $\bm{\alpha}_k$ is an eigenvector of $\Sigma$ corresponding to the $k$th largest eigenvalue $\lambda_k$
\vfill
If, additionally, $\|\bm{\alpha}_k\|=\bm{\alpha}_k^T\bm{\alpha}=1$, then $\lambda_k=\Var z_k$
\end{frame}


\begin{frame}{Why is that?}
Let us start with
\[
\bm{\alpha}_1^T\bx
\]
\vfill
We want maximum variance, where $\bm{\alpha}_1=(\alpha_{11},\ldots,\alpha_{1p})$, i.e.,
\[
\bm{\alpha}_1^T\bx = \sum_{j=1}^p\alpha_{1j}x_j
\]
with the constraint that $\|\bm{\alpha}_1\|=1$
\vfill
We have
\[
\Var \bm{\alpha}_1^T\bx
=\bm{\alpha}_1^T\Sigma\bm{\alpha}_1
\]
\end{frame}


\begin{frame}{Computation of $\Var(\bm{\alpha}_1^T\bx)$}
\begin{align*}
\Var(\bm{\alpha}_1^T\bx) 
&= E\left[\left(\bm{\alpha}_1^T\bx - E\left(\bm{\alpha}_1^T\bx\right)\right)^2\right] \\
&= E\left[\left(\bm{\alpha}_1^T\bx - \bm{\alpha}_1^T E(\bx)\right)^2\right] \\
&= E\left[\left(\bm{\alpha}_1^T(\bx - E(\bx))\right)^2\right] \\
&= E\left[\left(\bm{\alpha}_1^T(\bx - E(\bx))\right)\left(\bm{\alpha}_1^T(\bx - E(\bx))\right)^T\right] \\
&= E\left[\bm{\alpha}_1^T(\bx - E(\bx))(\bx - E(\bx))^T\bm{\alpha}_1\right] \\
&= \bm{\alpha}_1^T E\left[(\bx - E(\bx))(\bx - E(\bx))^T\right]\bm{\alpha}_1 \\
&= \bm{\alpha}_1^T\Sigma\bm{\alpha}_1
\end{align*}
\end{frame}


\begin{frame}{Objective}
We want to maximise $\Var \bm{\alpha}_1^T\bx$, i.e.,
\[
\bm{\alpha}_1^T\Sigma\bm{\alpha}_1
\]
under the constraint that $\|\bm{\alpha}_1\|=1$
\vfill
$\implies$ use \textbf{Lagrange multipliers}
\end{frame}


\begin{frame}{Maximisation using Lagrange multipliers}
\framesubtitle{(Recall from the mathematical preliminaries slides)}
As presented in Unit~02 (\emph{Mathematical preliminaries}), the method of \textbf{Lagrange multipliers} establishes that to find the extremum of a function $f(x_1,\ldots,x_n)$ under the constraint $g(x_1,\ldots,x_n)=k$:
\begin{enumerate}
\item Solve
\begin{align*}
\nabla f(x_1,\ldots,x_n) &= \lambda\nabla g(x_1,\ldots,x_n) \\
g(x_1,\ldots,x_n) &= k
\end{align*}
where $\nabla=\left(\frac{\partial}{\partial x_1},\ldots,\frac{\partial}{\partial x_n}\right)$ is the \textbf{gradient operator}
\item Plug all solutions into $f(x_1,\ldots,x_n)$ and find maximum values (provided values exist and $\nabla g\neq \b0$ there)
\end{enumerate}
\vfill
$\lambda\in\IR$ is the \textbf{Lagrange multiplier}
\vfill
Geometrically, this corresponds to the tangency condition: gradient vectors are \textbf{parallel} ($\nabla f = \lambda \nabla g$)
\end{frame}


\begin{frame}{Bear with me..}
\framesubtitle{(You may experience a brief period of discomfort)}
$\bm{\alpha}_1^T\Sigma\bm{\alpha}_1$ and $\|\bm{\alpha}_1\|^2=\bm{\alpha}_1^T\bm{\alpha_1}$ are functions of $\bm{\alpha}_1=(\alpha_{11},\ldots,\alpha_{1p})$
\vfill
In the notation of the previous slide, we want the max of 
\[
f(\alpha_{11},\ldots,\alpha_{1p}) := \bm{\alpha}_1^T\Sigma\bm{\alpha}_1
\]
under the constraint that
\[
g(\alpha_{11},\ldots,\alpha_{1p}) := \bm{\alpha}_1^T\bm{\alpha_1} = 1
\]
and with gradient operator
\[
\nabla = \left(
\frac{\partial}{\partial \alpha_{11}},
\ldots,
\frac{\partial}{\partial \alpha_{1p}}
\right)
\]
\end{frame}


\begin{frame}{Effect of $\nabla$ on $g$}
$g$ is easiest to see:
\begin{align*}
\nabla g(\alpha_{11},\ldots,\alpha_{1p}) &=
\left(
\frac{\partial}{\partial \alpha_{11}},
\ldots,
\frac{\partial}{\partial \alpha_{1p}}
\right) (\alpha_{11},\ldots,\alpha_{1p}) 
\begin{pmatrix}
\alpha_{11}\\ \vdots\\ \alpha_{1p}
\end{pmatrix} \\
&= \left(
\frac{\partial}{\partial \alpha_{11}},
\ldots,
\frac{\partial}{\partial \alpha_{1p}}
\right) 
\left(
\alpha_{11}^2+\cdots+\alpha_{1p}^2
\right) \\
&= \left(2\alpha_{11},\ldots,2\alpha_{1p}\right)\\
&= 2\bm{\alpha}_1
\end{align*}
\vfill
(And that's a general result: $\nabla\|\bx\|_2^2=2\bx$ with $\|\cdot\|_2$ the Euclidean norm)
\end{frame}

\begin{frame}{Effect of $\nabla$ on $f$}
Expand (write $\Sigma=[s_{ij}]$ and do not exploit symmetry)
\begin{align*}
\bm{\alpha}_1^T\Sigma\bm{\alpha}_1 &=
\left(\alpha_{11},\ldots,\alpha_{1p}\right)
\begin{pmatrix}
s_{11} & s_{12} & \cdots & s_{1p} \\
s_{21} & s_{22} & \cdots & s_{2p} \\
\vdots & \vdots & & \vdots \\
s_{p1} & s_{p2} & & s_{pp}
\end{pmatrix}
\begin{pmatrix}
\alpha_{11} \\ \alpha_{12} \\ \vdots \\ \alpha_{1p}
\end{pmatrix} \\
&=
\left(\alpha_{11},\ldots,\alpha_{1p}\right)
\begin{pmatrix}
s_{11}\alpha_{11}+s_{12}\alpha_{12}+\cdots+s_{1p}\alpha_{1p} \\
s_{21}\alpha_{11}+s_{22}\alpha_{12}+\cdots+s_{2p}\alpha_{1p} \\
\vdots \\
s_{p1}\alpha_{11}+s_{p2}\alpha_{12}+\cdots+s_{pp}\alpha_{1p}
\end{pmatrix} \\
&=
(s_{11}\alpha_{11}+s_{12}\alpha_{12}+\cdots+s_{1p}\alpha_{1p})\alpha_{11} \\
&\quad +
(s_{21}\alpha_{11}+s_{22}\alpha_{12}+\cdots+s_{2p}\alpha_{1p})\alpha_{12} \\
&\quad\;\;\vdots \\
&\quad +
(s_{p1}\alpha_{11}+s_{p2}\alpha_{12}+\cdots+s_{pp}\alpha_{1p})\alpha_{1p}
\end{align*}
\end{frame}

\begin{frame}
We have
\begin{align*}
\bm{\alpha}_1^T\Sigma\bm{\alpha}_1 &=
(s_{11}\alpha_{11}+s_{12}\alpha_{12}+\cdots+s_{1p}\alpha_{1p})\alpha_{11} \\
&\quad +
(s_{21}\alpha_{11}+s_{22}\alpha_{12}+\cdots+s_{2p}\alpha_{1p})\alpha_{12} \\
&\quad\;\;\vdots \\
&\quad +
(s_{p1}\alpha_{11}+s_{p2}\alpha_{12}+\cdots+s_{pp}\alpha_{1p})\alpha_{1p} 
\end{align*}
\begin{align*}
\implies\frac{\partial}{\partial \alpha_{11}}
\bm{\alpha}_1^T\Sigma\bm{\alpha}_1  
&= 
(s_{11}\alpha_{11}+s_{12}\alpha_{12}+\cdots+s_{1p}\alpha_{1p})+s_{11}\alpha_{11} \\
&\quad + s_{21}\alpha_{12} +\cdots + s_{p1}\alpha_{1p} \\
&= s_{11}\alpha_{11}+s_{12}\alpha_{12}+\cdots+s_{1p}\alpha_{1p} \\
&\quad+
s_{11}\alpha_{11}+s_{21}\alpha_{12}+\cdots+s_{p1}\alpha_{1p} \\
&= 2(s_{11}\alpha_{11}+s_{12}\alpha_{12}+\cdots+s_{1p}\alpha_{1p})
\end{align*}
(last equality stems from symmetry of $\Sigma$)
\end{frame}

\begin{frame}
In general, for $i=1,\ldots,p$,
\begin{align*}
\frac{\partial}{\partial \alpha_{1i}}
\bm{\alpha}_1^T\Sigma\bm{\alpha}_1  
&= s_{i1}\alpha_{11}+s_{i2}\alpha_{12}+\cdots+s_{ip}\alpha_{1p}\\
&\quad+s_{i1}\alpha_{11}+s_{2i}\alpha_{12}+\cdots+s_{pi}\alpha_{1p} \\
&= 2(s_{i1}\alpha_{11}+s_{i2}\alpha_{12}+\cdots+s_{ip}\alpha_{1p})
\end{align*}
(because of symmetry of $\Sigma$)
\vfill
As a consequence,
\[
\nabla \bm{\alpha}_1^T\Sigma\bm{\alpha}_1
=2\Sigma\bm{\alpha}_1
\]
\end{frame}

\begin{frame}
So solving
\[
\nabla f(x_1,\ldots,x_n) = \lambda\nabla g(x_1,\ldots,x_n) 
\]
means solving
\[
2\Sigma\bm{\alpha}_1 = \lambda 2\bm{\alpha}_1 
\]
i.e.,
\[
\Sigma\bm{\alpha}_1 = \lambda\bm{\alpha}_1 
\]
\vfill
$\implies$
$(\lambda,\bm{\alpha}_1)$ eigenpair of $\Sigma$, with $\bm{\alpha}_1$ having unit length
\end{frame}


\begin{frame}{Picking the right eigenvalue}
$(\lambda,\bm{\alpha}_1)$ eigenpair of $\Sigma$, with $\bm{\alpha}_1$ having unit length
\vfill
But which $\lambda$ to choose?
\vfill
Recall that we want $\Var \bm{\alpha}_1^T\bx=\bm{\alpha}_1^T\Sigma\bm{\alpha}_1$ maximal
\vfill
We have
\[
\Var \bm{\alpha}_1^T\bx 
= \bm{\alpha}_1^T\Sigma\bm{\alpha}_1 
= \bm{\alpha}_1^T(\Sigma\bm{\alpha}_1) 
= \bm{\alpha}_1^T(\lambda\bm{\alpha}_1) 
= \lambda(\bm{\alpha}_1^T\bm{\alpha}_1) = \lambda
\]
\vfill
$\implies$ we pick $\lambda=\lambda_1$, the largest eigenvalue (covariance matrix symmetric so eigenvalues real)
\end{frame}


\begin{frame}{What we have this far..}
The first principal component is $\bm{\alpha}_1^T\bx$ and has variance $\lambda_1$, where $\lambda_1$ the largest eigenvalue of $\Sigma$ and $\bm{\alpha}_1$ an associated eigenvector with $\|\bm{\alpha}_1\|=1$
\vfill
We want the second principal component to be \emph{uncorrelated} with $\bm{\alpha}_1^T\bx$ and to have maximum variance $\Var \bm{\alpha}_2^T\bx=\bm{\alpha}_2^T\Sigma\bm{\alpha}_2$, under the constraint that $\|\bm{\alpha}_2\|=1$
\vfill
$\bm{\alpha}_2^T\bx$ uncorrelated to $\bm{\alpha}_1^T\bx$ if $\cov(\bm{\alpha}_1^T\bx,\bm{\alpha}_2^T\bx)=0$
\end{frame}

\begin{frame}
We have
\begin{align*}
\cov(\bm{\alpha}_1^T\bx,\bm{\alpha}_2^T\bx) &= 
\bm{\alpha}_1^T\Sigma\bm{\alpha}_2 \\
&= \bm{\alpha}_2^T\Sigma^T\bm{\alpha}_1 \\
&= \bm{\alpha}_2^T\Sigma\bm{\alpha}_1 \quad\textrm{[$\Sigma$ symmetric]} \\
&= \bm{\alpha}_2^T(\lambda_1\bm{\alpha}_1) \\
&= \lambda \bm{\alpha}_2^T\bm{\alpha}_1
\end{align*}
\vfill
So $\bm{\alpha}_2^T\bx$ uncorrelated to $\bm{\alpha}_1^T\bx$ if $\bm{\alpha}_1\perp\bm{\alpha}_2$
\vfill
This is beginning to sound a lot like Gram-Schmidt, no?
\end{frame}

\begin{frame}{In short}
Take whatever covariance matrix is available to you (known $\Sigma$ or sample $S_X$) -- assume sample from now on for simplicity
\vfill
For $i=1,\ldots,p$, the $i$th principal component is
\[
z_i = \bv_i^T\bx
\]
where $\bv_i$ eigenvector of $S_X$ associated to the $i$th largest eigenvalue $\lambda_i$
\vfill
If $\bv_i$ is normalised, then $\lambda_i=\Var z_i$
\end{frame}


\begin{frame}{A 3D example -- Sample covariance matrix}
Suppose we have a dataset with $p=3$ variables, $\bx = (x_1, x_2, x_3)^T$, centred so that sample means are zero
\vfill
After computing all sample variances and covariances, we obtain the sample covariance matrix:
\[
S_X = \begin{pmatrix}
15 & 10 & 4 \\
10 & 14 & 6 \\
4 & 6 & 7
\end{pmatrix}
\]
\vfill
\begin{itemize}
    \item Sample variances (diagonal): $\Var(x_1) = 15$, $\Var(x_2) = 14$, $\Var(x_3) = 7$
    \item Sample covariances (off-diagonal): all variables are positively correlated ($\cov(x_1,x_2)=10$, $\cov(x_1,x_3)=4$, $\cov(x_2,x_3)=6$)
    \item Total sample variance in the dataset is the trace:
    \[
    \tr(S_X) = 15 + 14 + 7 = 36
    \]
\end{itemize}
\end{frame}


\begin{frame}{A 3D example -- Eigenvalues}
To compute the principal components, we find the eigenvalues of $S_X$ by solving $\det(S_X - \lambda I) = 0$:
\[
\det\begin{pmatrix}
15 - \lambda & 10 & 4 \\
10 & 14 - \lambda & 6 \\
4 & 6 & 7 - \lambda
\end{pmatrix} = 0
\]
\vfill
This yields three distinct positive eigenvalues:
\[
\lambda_1 = 27, \qquad \lambda_2 = 6, \qquad \lambda_3 = 3
\]
\vfill
Note that:
\[
\lambda_1 + \lambda_2 + \lambda_3 = 27 + 6 + 3 = 36 = \tr(S_X)
\]
As expected from theory, the total variance is preserved!
\end{frame}


\begin{frame}{A 3D example -- Variance explained}
The variance of each principal component $z_i = \bv_i^T\bx$ is given by $\Var z_i = \lambda_i$ (for normalized $\bv_i$)
\vfill
Proportion of total variance explained by each principal component:
\begin{itemize}
    \item \textbf{1st PC ($z_1$)}: 
    \[
    \frac{\lambda_1}{\tr(S_X)} = \frac{27}{36} = \frac{3}{4} = 75\%
    \]
    \item \textbf{2nd PC ($z_2$)}: 
    \[
    \frac{\lambda_2}{\tr(S_X)} = \frac{6}{36} = \frac{1}{6} \approx 16.67\%
    \]
    \item \textbf{3rd PC ($z_3$)}: 
    \[
    \frac{\lambda_3}{\tr(S_X)} = \frac{3}{36} = \frac{1}{12} \approx 8.33\%
    \]
\end{itemize}
\vfill
The first two principal components together account for $\frac{27+6}{36} = \frac{33}{36} \approx 91.67\%$ of the total variance
\end{frame}


\begin{frame}{A 3D example -- Eigenvectors}
Next, we find the associated normalized eigenvectors by solving $(S_X - \lambda_i I)\bv_i = \b0$ subject to $\|\bv_i\| = 1$:
\vfill
\begin{itemize}
    \item For $\lambda_1 = 27$:
    \[
    \bv_1 = \frac{1}{3}\begin{pmatrix} 2 \\ 2 \\ 1 \end{pmatrix} = \begin{pmatrix} 2/3 \\ 2/3 \\ 1/3 \end{pmatrix}, \qquad \|\bv_1\|^2 = \frac{4+4+1}{9} = 1
    \]
    \item For $\lambda_2 = 6$:
    \[
    \bv_2 = \frac{1}{3}\begin{pmatrix} -2 \\ 1 \\ 2 \end{pmatrix} = \begin{pmatrix} -2/3 \\ 1/3 \\ 2/3 \end{pmatrix}, \qquad \|\bv_2\|^2 = \frac{4+1+4}{9} = 1
    \]
    \item For $\lambda_3 = 3$:
    \[
    \bv_3 = \frac{1}{3}\begin{pmatrix} 1 \\ -2 \\ 2 \end{pmatrix} = \begin{pmatrix} 1/3 \\ -2/3 \\ 2/3 \end{pmatrix}, \qquad \|\bv_3\|^2 = \frac{1+4+4}{9} = 1
    \]
\end{itemize}
\end{frame}


\begin{frame}{A 3D example -- Orthonormality check}
Since $S_X$ is symmetric with distinct eigenvalues, the eigenvectors are mutually orthogonal:
\begin{align*}
\bv_1^T\bv_2 &= \frac{1}{9}\big(2(-2) + 2(1) + 1(2)\big) = \frac{-4 + 2 + 2}{9} = 0 \\
\bv_1^T\bv_3 &= \frac{1}{9}\big(2(1) + 2(-2) + 1(2)\big) = \frac{2 - 4 + 2}{9} = 0 \\
\bv_2^T\bv_3 &= \frac{1}{9}\big(-2(1) + 1(-2) + 2(2)\big) = \frac{-2 - 2 + 4}{9} = 0
\end{align*}
\vfill
Hence, $V = \begin{pmatrix} \bv_1 & \bv_2 & \bv_3 \end{pmatrix}$ is an orthogonal matrix ($V^TV = I$), forming an orthonormal basis of $\IR^3$
\end{frame}


\begin{frame}{A 3D example -- The principal components}
The principal components are the linear combinations $z_i = \bv_i^T\bx$:
\begin{align*}
z_1 &= \bv_1^T\bx = \frac{2}{3}x_1 + \frac{2}{3}x_2 + \frac{1}{3}x_3 \\
z_2 &= \bv_2^T\bx = -\frac{2}{3}x_1 + \frac{1}{3}x_2 + \frac{2}{3}x_3 \\
z_3 &= \bv_3^T\bx = \frac{1}{3}x_1 - \frac{2}{3}x_2 + \frac{2}{3}x_3
\end{align*}
\vfill
Properties of these new variables:
\begin{itemize}
    \item $\Var z_1 = \lambda_1 = 27$, $\Var z_2 = \lambda_2 = 6$, $\Var z_3 = \lambda_3 = 3$
    \item $\cov(z_i, z_j) = 0$ for all $i \neq j$ (uncorrelated)
    \item In the new basis, the covariance matrix is diagonal:
    \[
    S_Z = V^T S_X V = \diag(27, 6, 3)
    \]
\end{itemize}
\end{frame}


\begin{frame}{A 3D example -- Dimensionality reduction}
Suppose we want to reduce the dimension of the data from $3$ to $2$:
\vfill
\begin{itemize}
    \item Keep only the first two components, $z_1$ and $z_2$
    \item $z_1$ and $z_2$ together preserve $91.67\%$ of the total variance
    \item Each 3D data point $\bx=(x_1,x_2,x_3)^T$ is mapped to its 2D coordinates $(z_1,z_2)^T$
    \item Approximation of $\bx$ using only 2 components:
    \[
    \bx \approx z_1 \bv_1 + z_2 \bv_2
    \]
    \item Discarding $z_3$ loses only $8.33\%$ of the variance
\end{itemize}
\end{frame}


\begin{frame}[fragile]{A 3D example -- Verifying in R}
We can confirm our calculations directly in R:
\begin{knitrout}
\definecolor{shadecolor}{rgb}{0.969, 0.969, 0.969}\color{fgcolor}\begin{kframe}
\begin{alltt}
\hldef{S} \hlkwb{=} \hlkwd{matrix}\hldef{(}\hlkwd{c}\hldef{(}\hlnum{15}\hldef{,} \hlnum{10}\hldef{,} \hlnum{4}\hldef{,} \hlnum{10}\hldef{,} \hlnum{14}\hldef{,} \hlnum{6}\hldef{,} \hlnum{4}\hldef{,} \hlnum{6}\hldef{,} \hlnum{7}\hldef{),} \hlkwc{nrow} \hldef{=} \hlnum{3}\hldef{,} \hlkwc{byrow} \hldef{=} \hlnum{TRUE}\hldef{)}
\hldef{ev} \hlkwb{=} \hlkwd{eigen}\hldef{(S)}
\hldef{ev}\hlopt{$}\hldef{values}
\end{alltt}
\begin{verbatim}
## [1] 27  6  3
\end{verbatim}
\begin{alltt}
\hlcom{# Proportion and cumulative proportion of variance explained}
\hldef{ev}\hlopt{$}\hldef{values} \hlopt{/} \hlkwd{sum}\hldef{(ev}\hlopt{$}\hldef{values)}
\end{alltt}
\begin{verbatim}
## [1] 0.75000000 0.16666667 0.08333333
\end{verbatim}
\begin{alltt}
\hlkwd{cumsum}\hldef{(ev}\hlopt{$}\hldef{values)} \hlopt{/} \hlkwd{sum}\hldef{(ev}\hlopt{$}\hldef{values)}
\end{alltt}
\begin{verbatim}
## [1] 0.7500000 0.9166667 1.0000000
\end{verbatim}
\end{kframe}
\end{knitrout}
\end{frame}





%%%%%%%%%%%%%%%%%%%%%
%%%%%%%%%%%%%%%%%%%%%
\SSsection{Hockey players}



\begin{frame}[fragile]\frametitle{Covariance of height and weight}
The function \code{cov} returns the sample covariance of two variables
\vfill
Recall that \code{cov} yields the same result whether the data is centred or uncentred:
\vfill
\begin{knitrout}
\definecolor{shadecolor}{rgb}{0.969, 0.969, 0.969}\color{fgcolor}\begin{kframe}
\begin{alltt}
\hlkwd{cov}\hldef{(data}\hlopt{$}\hldef{height, data}\hlopt{$}\hldef{weight)}
\end{alltt}
\begin{verbatim}
## [1] 26.63506
\end{verbatim}
\begin{alltt}
\hlkwd{cov}\hldef{(data}\hlopt{$}\hldef{height.c, data}\hlopt{$}\hldef{weight.c)}
\end{alltt}
\begin{verbatim}
## [1] 26.63506
\end{verbatim}
\end{kframe}
\end{knitrout}
\vfill
As seen in the scatter plot, there is a positive linear relationship between height and weight
\end{frame}


\begin{frame}[fragile]\frametitle{Sample covariance matrix}
Let us form the $n\times 2$ centred data matrix $\widetilde{X}$ and compute the sample covariance matrix $S = \frac{1}{n-1}\widetilde{X}^T\widetilde{X}$
\vfill
\begin{knitrout}
\definecolor{shadecolor}{rgb}{0.969, 0.969, 0.969}\color{fgcolor}\begin{kframe}
\begin{alltt}
\hldef{X} \hlkwb{=} \hlkwd{as.matrix}\hldef{(data[,} \hlkwd{c}\hldef{(}\hlsng{"height.c"}\hldef{,} \hlsng{"weight.c"}\hldef{)])}
\hldef{S} \hlkwb{=} \hlnum{1} \hlopt{/} \hldef{(}\hlkwd{dim}\hldef{(X)[}\hlnum{1}\hldef{]} \hlopt{-} \hlnum{1}\hldef{)} \hlopt{*} \hlkwd{t}\hldef{(X)} \hlopt{%*%} \hldef{X}
\hldef{S}
\end{alltt}
\begin{verbatim}
##          height.c weight.c
## height.c 29.66176 26.63506
## weight.c 26.63506 47.81112
\end{verbatim}
\end{kframe}
\end{knitrout}
\vfill
The off-diagonal entries match $\cov(X_1, X_2)$
\end{frame}


\begin{frame}[fragile]\frametitle{Sample covariance matrix -- Variances}
Let us check that the diagonal entries indeed match the individual sample variances:
\vfill
\begin{knitrout}
\definecolor{shadecolor}{rgb}{0.969, 0.969, 0.969}\color{fgcolor}\begin{kframe}
\begin{alltt}
\hlkwd{var}\hldef{(X[,} \hlnum{1}\hldef{])}
\end{alltt}
\begin{verbatim}
## [1] 29.66176
\end{verbatim}
\begin{alltt}
\hlkwd{var}\hldef{(X[,} \hlnum{2}\hldef{])}
\end{alltt}
\begin{verbatim}
## [1] 47.81112
\end{verbatim}
\end{kframe}
\end{knitrout}
\vfill
Both match perfectly
\end{frame}


\begin{frame}[fragile]\frametitle{Computing the principal components}
We compute the eigenvalues and eigenvectors of $S$ using \code{eigen}
\vfill
\begin{knitrout}
\definecolor{shadecolor}{rgb}{0.969, 0.969, 0.969}\color{fgcolor}\begin{kframe}
\begin{alltt}
\hldef{ev} \hlkwb{=} \hlkwd{eigen}\hldef{(S)}
\hldef{ev}
\end{alltt}
\begin{verbatim}
## eigen() decomposition
## $values
## [1] 66.87496 10.59793
## 
## $vectors
##           [,1]       [,2]
## [1,] 0.5820222 -0.8131729
## [2,] 0.8131729  0.5820222
\end{verbatim}
\end{kframe}
\end{knitrout}
\vfill
\code{eigen} returns eigenvalues in decreasing order and orthonormal eigenvectors
\end{frame}


\begin{frame}\frametitle{Variance explained}
The total variance is
\[
\tr(S) = \lambda_1 + \lambda_2 = 66.87 + 10.6 = 77.47
\]
\vfill
The first principal component explains
\[
\frac{\lambda_1}{\lambda_1+\lambda_2} = \frac{66.87}{77.47} \approx 86.3\%
\]
of the total variance in the height and weight data
\vfill
The second principal component accounts for the remaining $13.7\%$
\end{frame}


\begin{frame}[fragile]\frametitle{Plotting the first principal component}
We superimpose the line spanned by the first eigenvector $\bv_1$ on the centred data
\vfill
The line passes through $(0,0)$ and has slope $b = v_{21}/v_{11}$
\vfill
To avoid visual distortion of angles, we use a square plot (\code{pty = "s"}, \code{asp = 1}) with equal ranges on both axes
\vfill
\begin{knitrout}
\definecolor{shadecolor}{rgb}{0.969, 0.969, 0.969}\color{fgcolor}\begin{kframe}
\begin{alltt}
\hlkwd{par}\hldef{(}\hlkwc{pty} \hldef{=} \hlsng{"s"}\hldef{)}
\hldef{lims} \hlkwb{=} \hlkwd{range}\hldef{(}\hlkwd{c}\hldef{(data}\hlopt{$}\hldef{height.c, data}\hlopt{$}\hldef{weight.c))}
\hlkwd{plot}\hldef{(data}\hlopt{$}\hldef{height.c, data}\hlopt{$}\hldef{weight.c,}
     \hlkwc{xlim} \hldef{= lims,} \hlkwc{ylim} \hldef{= lims,} \hlkwc{asp} \hldef{=} \hlnum{1}\hldef{,}
     \hlkwc{pch} \hldef{=} \hlnum{19}\hldef{,} \hlkwc{col} \hldef{=} \hlsng{"dodgerblue4"}\hldef{,}
     \hlkwc{main} \hldef{=} \hlsng{"IIHF players 2001-2016 (with first component)"}\hldef{,}
     \hlkwc{xlab} \hldef{=} \hlsng{"Centred height (cm)"}\hldef{,} \hlkwc{ylab} \hldef{=} \hlsng{"Centred weight (kg)"}\hldef{)}
\hlkwd{abline}\hldef{(}\hlkwc{a} \hldef{=} \hlnum{0}\hldef{,} \hlkwc{b} \hldef{= ev}\hlopt{$}\hldef{vectors[}\hlnum{2}\hldef{,} \hlnum{1}\hldef{]} \hlopt{/} \hldef{ev}\hlopt{$}\hldef{vectors[}\hlnum{1}\hldef{,} \hlnum{1}\hldef{],}
       \hlkwc{col} \hldef{=} \hlsng{"red"}\hldef{,} \hlkwc{lwd} \hldef{=} \hlnum{3}\hldef{)}
\end{alltt}
\end{kframe}
\end{knitrout}
\end{frame}

\maxFrameImage{FIGS/U05-plot-hockey-centred-evector-1.pdf}


\begin{frame}[fragile]\frametitle{First principal component vs.\ least squares}
How does the first principal component line compare with ordinary least squares (OLS) regression?
\vfill
Fit a linear model with zero intercept (since data is centred):
\vfill
\begin{knitrout}
\definecolor{shadecolor}{rgb}{0.969, 0.969, 0.969}\color{fgcolor}\begin{kframe}
\begin{alltt}
\hldef{z} \hlkwb{=} \hlkwd{lm}\hldef{(weight.c} \hlopt{~} \hlnum{0} \hlopt{+} \hldef{height.c,} \hlkwc{data} \hldef{= data)}
\hlkwd{coef}\hldef{(z)}
\end{alltt}
\begin{verbatim}
##  height.c 
## 0.8979593
\end{verbatim}
\end{kframe}
\end{knitrout}
\vfill

\end{frame}

\maxFrameImage{FIGS/U05-plot-hockey-lm-compare-1.pdf}


\begin{frame}{Why are the two lines different?}
\textbf{Ordinary least squares (OLS)}:
\begin{itemize}
\item Treats \code{height.c} as an exact predictor
\item Minimises the sum of squared \textbf{vertical} residuals (errors parallel to the $y$-axis)
\end{itemize}
\vfill
\textbf{Principal component analysis (PCA)}:
\begin{itemize}
\item Treats both variables \textbf{symmetrically}
\item Maximises the variance of the orthogonal projection
\item Equivalently, minimises the sum of squared \textbf{orthogonal (perpendicular)} distances from points to the line
\end{itemize}
\end{frame}


\begin{frame}\frametitle{Rotating the data}
Let us rotate the data so that the first principal component becomes the horizontal axis
\vfill
A counter-clockwise rotation by an angle $\phi$ in $\IR^2$ is represented by
\[
R_\phi = \begin{pmatrix} \cos\phi & -\sin\phi \\ \sin\phi & \cos\phi \end{pmatrix}
\]
\vfill
The first eigenvector makes an angle $\theta > 0$ with the positive $x$-axis, $\tan\theta = v_{21}/v_{11}$
\vfill
To align it with the horizontal axis, we rotate \textbf{clockwise} by $\theta$, i.e., by $\phi = -\theta$
\end{frame}


\begin{frame}[fragile]\frametitle{Computing the rotation matrix}
\begin{knitrout}
\definecolor{shadecolor}{rgb}{0.969, 0.969, 0.969}\color{fgcolor}\begin{kframe}
\begin{alltt}
\hldef{theta} \hlkwb{=} \hlkwd{atan}\hldef{(ev}\hlopt{$}\hldef{vectors[}\hlnum{2}\hldef{,} \hlnum{1}\hldef{]} \hlopt{/} \hldef{ev}\hlopt{$}\hldef{vectors[}\hlnum{1}\hldef{,} \hlnum{1}\hldef{])}
\hldef{theta}
\end{alltt}
\begin{verbatim}
## [1] 0.949583
\end{verbatim}
\begin{alltt}
\hldef{R_minus_theta} \hlkwb{=} \hlkwd{matrix}\hldef{(}\hlkwd{c}\hldef{(}\hlkwd{cos}\hldef{(}\hlopt{-}\hldef{theta),} \hlopt{-}\hlkwd{sin}\hldef{(}\hlopt{-}\hldef{theta),}
                         \hlkwd{sin}\hldef{(}\hlopt{-}\hldef{theta),}  \hlkwd{cos}\hldef{(}\hlopt{-}\hldef{theta)),}
                       \hlkwc{nr} \hldef{=} \hlnum{2}\hldef{,} \hlkwc{byrow} \hldef{=} \hlnum{TRUE}\hldef{)}
\hldef{R_minus_theta}
\end{alltt}
\begin{verbatim}
##            [,1]      [,2]
## [1,]  0.5820222 0.8131729
## [2,] -0.8131729 0.5820222
\end{verbatim}
\end{kframe}
\end{knitrout}
\end{frame}


\begin{frame}[fragile]\frametitle{Rotating the data}
Apply the rotation matrix to each observation:
\vfill
\begin{knitrout}
\definecolor{shadecolor}{rgb}{0.969, 0.969, 0.969}\color{fgcolor}\begin{kframe}
\begin{alltt}
\hldef{rot_pts} \hlkwb{=} \hldef{X} \hlopt{%*%} \hlkwd{t}\hldef{(R_minus_theta)}
\hldef{data}\hlopt{$}\hldef{height.c_r} \hlkwb{=} \hldef{rot_pts[,} \hlnum{1}\hldef{]}
\hldef{data}\hlopt{$}\hldef{weight.c_r} \hlkwb{=} \hldef{rot_pts[,} \hlnum{2}\hldef{]}
\hlkwd{plot}\hldef{(data}\hlopt{$}\hldef{height.c_r, data}\hlopt{$}\hldef{weight.c_r,}
     \hlkwc{xlim} \hldef{= lims,} \hlkwc{ylim} \hldef{= lims,} \hlkwc{asp} \hldef{=} \hlnum{1}\hldef{,}
     \hlkwc{pch} \hldef{=} \hlnum{19}\hldef{,} \hlkwc{col} \hldef{=} \hlsng{"dodgerblue4"}\hldef{,}
     \hlkwc{xlab} \hldef{=} \hlsng{"Rotated x-axis (1st component)"}\hldef{,}
     \hlkwc{ylab} \hldef{=} \hlsng{"Rotated y-axis"}\hldef{,}
     \hlkwc{main} \hldef{=} \hlsng{"IIHF players 2001-2016 (rotated to first component)"}\hldef{)}
\hlkwd{abline}\hldef{(}\hlkwc{h} \hldef{=} \hlnum{0}\hldef{,} \hlkwc{col} \hldef{=} \hlsng{"red"}\hldef{,} \hlkwc{lwd} \hldef{=} \hlnum{2}\hldef{)}
\end{alltt}
\end{kframe}
\end{knitrout}
\end{frame}

\maxFrameImage{FIGS/U05-plot-hockey-centred-rotated-1.pdf}


\begin{frame}[fragile]\frametitle{First and second principal components}
Plot both eigenvectors on the centred data. Since $S$ is symmetric, the eigenvectors are orthogonal ($90^\circ$ angle):
\vfill
\begin{knitrout}
\definecolor{shadecolor}{rgb}{0.969, 0.969, 0.969}\color{fgcolor}\begin{kframe}
\begin{alltt}
\hlkwd{plot}\hldef{(data}\hlopt{$}\hldef{height.c, data}\hlopt{$}\hldef{weight.c,}
     \hlkwc{xlim} \hldef{= lims,} \hlkwc{ylim} \hldef{= lims,} \hlkwc{asp} \hldef{=} \hlnum{1}\hldef{,}
     \hlkwc{pch} \hldef{=} \hlnum{19}\hldef{,} \hlkwc{col} \hldef{=} \hlsng{"dodgerblue4"}\hldef{,}
     \hlkwc{main} \hldef{=} \hlsng{"IIHF players 2001-2016 (1st and 2nd eigenvectors)"}\hldef{,}
     \hlkwc{xlab} \hldef{=} \hlsng{"Centred height (cm)"}\hldef{,} \hlkwc{ylab} \hldef{=} \hlsng{"Centred weight (kg)"}\hldef{)}
\hlkwd{abline}\hldef{(}\hlkwc{a} \hldef{=} \hlnum{0}\hldef{,} \hlkwc{b} \hldef{= ev}\hlopt{$}\hldef{vectors[}\hlnum{2}\hldef{,} \hlnum{1}\hldef{]} \hlopt{/} \hldef{ev}\hlopt{$}\hldef{vectors[}\hlnum{1}\hldef{,} \hlnum{1}\hldef{],}
       \hlkwc{col} \hldef{=} \hlsng{"red"}\hldef{,} \hlkwc{lwd} \hldef{=} \hlnum{3}\hldef{)}
\hlkwd{abline}\hldef{(}\hlkwc{a} \hldef{=} \hlnum{0}\hldef{,} \hlkwc{b} \hldef{= ev}\hlopt{$}\hldef{vectors[}\hlnum{2}\hldef{,} \hlnum{2}\hldef{]} \hlopt{/} \hldef{ev}\hlopt{$}\hldef{vectors[}\hlnum{1}\hldef{,} \hlnum{2}\hldef{],}
       \hlkwc{col} \hldef{=} \hlsng{"darkgreen"}\hldef{,} \hlkwc{lwd} \hldef{=} \hlnum{3}\hldef{)}
\end{alltt}
\end{kframe}
\end{knitrout}
\end{frame}

\maxFrameImage{FIGS/U05-plot-hockey-centred-2evectors-1.pdf}


\begin{frame}[fragile]\frametitle{Representing data in the new coordinate system}
Rather than 2D trigonometry, express data in the orthonormal eigenbasis using our general change-of-basis machinery
\vfill
Enforce orthonormality with \code{pracma::gramSchmidt} and compute the change-of-basis matrix using \code{pracma::rref}:
\vfill
\begin{knitrout}
\definecolor{shadecolor}{rgb}{0.969, 0.969, 0.969}\color{fgcolor}\begin{kframe}
\begin{alltt}
\hldef{GS} \hlkwb{=} \hldef{pracma}\hlopt{::}\hlkwd{gramSchmidt}\hldef{(}\hlkwc{A} \hldef{= ev}\hlopt{$}\hldef{vectors,} \hlkwc{tol} \hldef{=} \hlnum{1e-10}\hldef{)}
\hldef{A} \hlkwb{=} \hlkwd{cbind}\hldef{(GS}\hlopt{$}\hldef{Q,} \hlkwd{diag}\hldef{(}\hlnum{2}\hldef{))}
\hldef{pracma}\hlopt{::}\hlkwd{rref}\hldef{(A)}
\end{alltt}
\begin{verbatim}
##      [,1] [,2]       [,3]      [,4]
## [1,]    1    0  0.5820222 0.8131729
## [2,]    0    1 -0.8131729 0.5820222
\end{verbatim}
\end{kframe}
\end{knitrout}
\vfill
$P = Q^{-1} = Q^T$ is the change-of-basis matrix (right-hand block)
\end{frame}


\begin{frame}[fragile]\frametitle{Projecting into the new basis}
Multiplying $\widetilde{X}$ by $P^T$ gives coordinates $X_{\text{new}}$ in the new PC coordinate system:
\vfill
\begin{knitrout}
\definecolor{shadecolor}{rgb}{0.969, 0.969, 0.969}\color{fgcolor}\begin{kframe}
\begin{alltt}
\hldef{P} \hlkwb{=} \hldef{pracma}\hlopt{::}\hlkwd{rref}\hldef{(A)[,} \hlkwd{c}\hldef{(}\hlnum{3}\hldef{,} \hlnum{4}\hldef{)]}
\hldef{P}
\end{alltt}
\begin{verbatim}
##            [,1]      [,2]
## [1,]  0.5820222 0.8131729
## [2,] -0.8131729 0.5820222
\end{verbatim}
\begin{alltt}
\hldef{X.new} \hlkwb{=} \hldef{X} \hlopt{%*%} \hlkwd{t}\hldef{(P)}
\hlkwd{cov}\hldef{(X.new)}
\end{alltt}
\begin{verbatim}
##               [,1]          [,2]
## [1,]  6.687496e+01 -2.328994e-13
## [2,] -2.328994e-13  1.059793e+01
\end{verbatim}
\end{kframe}
\end{knitrout}
\end{frame}

\maxFrameImage{FIGS/U05-plot-hockey-centred-evector-3-1.pdf}


\begin{frame}[fragile]\frametitle{Diagonal covariance in the new basis}
As confirmed by \code{cov(X.new)}:
\vfill
\begin{knitrout}
\definecolor{shadecolor}{rgb}{0.969, 0.969, 0.969}\color{fgcolor}\begin{kframe}
\begin{alltt}
\hlkwd{cov}\hldef{(X.new)}
\end{alltt}
\begin{verbatim}
##               [,1]          [,2]
## [1,]  6.687496e+01 -2.328994e-13
## [2,] -2.328994e-13  1.059793e+01
\end{verbatim}
\end{kframe}
\end{knitrout}
\vfill
\begin{itemize}
\item The sample covariance matrix in the new coordinate system is \textbf{diagonal}
\item Diagonal entries equal $\lambda_1 = 66.87$ and $\lambda_2 = 10.6$
\item Off-diagonal entry is zero: the new coordinates are completely \textbf{uncorrelated}
\end{itemize}
\end{frame}


%%%%%%%%%%%%%%%%%%%%%
%%%%%%%%%%%%%%%%%%%%%
\SSsubsection{PCA with dedicated R packages}

\begin{frame}[fragile]\frametitle{Multivariate hockey dataset}
In practice, we use dedicated packages such as \code{prcomp}, \code{FactoMineR} and \code{ggbiplot}, especially with $> 2$ variables
\vfill
We retain 4 numerical attributes (\code{height}, \code{weight}, \code{age}, \code{bmi}) and player \code{position} (\code{Defence}, \code{Forward}, \code{Goalkeeper}):
\vfill

\begin{knitrout}
\definecolor{shadecolor}{rgb}{0.969, 0.969, 0.969}\color{fgcolor}\begin{kframe}
\begin{alltt}
\hldef{agg_num} \hlkwb{=} \hlkwd{aggregate}\hldef{(}\hlkwd{cbind}\hldef{(height, weight, age, bmi)} \hlopt{~} \hldef{name,}
                    \hlkwc{data} \hldef{= raw_data,} \hlkwc{FUN} \hldef{= mean)}
\hldef{agg_pos} \hlkwb{=} \hlkwd{aggregate}\hldef{(position} \hlopt{~} \hldef{name,} \hlkwc{data} \hldef{= raw_data,}
                    \hlkwc{FUN} \hldef{=} \hlkwa{function}\hldef{(}\hlkwc{x}\hldef{) x[}\hlnum{1}\hldef{])}
\hldef{hockey_multi} \hlkwb{=} \hlkwd{merge}\hldef{(agg_pos, agg_num,} \hlkwc{by} \hldef{=} \hlsng{"name"}\hldef{)}
\hldef{hockey_multi}\hlopt{$}\hldef{position} \hlkwb{=} \hlkwd{factor}\hldef{(}
  \hldef{hockey_multi}\hlopt{$}\hldef{position,}
  \hlkwc{levels} \hldef{=} \hlkwd{c}\hldef{(}\hlsng{"D"}\hldef{,} \hlsng{"F"}\hldef{,} \hlsng{"G"}\hldef{),}
  \hlkwc{labels} \hldef{=} \hlkwd{c}\hldef{(}\hlsng{"Defence"}\hldef{,} \hlsng{"Forward"}\hldef{,} \hlsng{"Goalkeeper"}\hldef{)}
\hldef{)}
\end{alltt}
\end{kframe}
\end{knitrout}
\end{frame}


\begin{frame}[fragile]\frametitle{Multivariate hockey dataset}
The aggregated multivariate dataset:
\vfill
\begin{knitrout}
\definecolor{shadecolor}{rgb}{0.969, 0.969, 0.969}\color{fgcolor}\begin{kframe}
\begin{alltt}
\hlkwd{head}\hldef{(hockey_multi,} \hlkwc{n} \hldef{=} \hlnum{4}\hldef{)}
\end{alltt}
\begin{verbatim}
##                 name position height weight      age      bmi
## 1     aagaard mikkel  Forward    180     81 20.53662 25.00000
## 2       aalto antti   Forward    187     96 27.15948 27.45289
## 3 aaltonen juhamatti  Forward    184     87 27.15606 25.69707
## 4  abdelkader justin  Forward    185     97 26.17933 28.34186
\end{verbatim}
\end{kframe}
\end{knitrout}
\end{frame}


\begin{frame}[fragile]\frametitle{PCA using \code{prcomp}}
Because variables are in different physical units (cm, kg, years, $\text{kg}/\text{m}^2$), we standardise to unit variance (\code{scale. = TRUE}):
\vfill
\begin{knitrout}
\definecolor{shadecolor}{rgb}{0.969, 0.969, 0.969}\color{fgcolor}\begin{kframe}
\begin{alltt}
\hldef{hockey.pca} \hlkwb{=} \hlkwd{prcomp}\hldef{(}
  \hldef{hockey_multi[,} \hlkwd{c}\hldef{(}\hlsng{"height"}\hldef{,} \hlsng{"weight"}\hldef{,} \hlsng{"age"}\hldef{,} \hlsng{"bmi"}\hldef{)],}
  \hlkwc{center} \hldef{=} \hlnum{TRUE}\hldef{,} \hlkwc{scale.} \hldef{=} \hlnum{TRUE}
\hldef{)}
\hlkwd{summary}\hldef{(hockey.pca)}
\end{alltt}
\begin{verbatim}
## Importance of components:
##                           PC1    PC2    PC3     PC4
## Standard deviation     1.3933 1.1044 0.9154 0.03081
## Proportion of Variance 0.4853 0.3049 0.2095 0.00024
## Cumulative Proportion  0.4853 0.7903 0.9998 1.00000
\end{verbatim}
\end{kframe}
\end{knitrout}
\end{frame}


\begin{frame}\frametitle{Importance of components}
\begin{itemize}
\item PC1 ($48.5\%$) and PC2 ($30.5\%$) explain $79.0\%$ of the total variance across the four standardised variables
\vfill
\item The first three components capture $99.98\%$ of the total variance
\vfill
\item PC4 accounts for almost zero variance ($0.02\%$):
\[
\text{BMI} = \frac{\text{weight}}{\text{height}^2}
\]
Around the sample mean, BMI is nearly a linear combination of \code{height} and \code{weight}, so the 4D point cloud is practically three-dimensional!
\end{itemize}
\end{frame}


\begin{frame}[fragile]\frametitle{PCA biplot with \code{ggbiplot}}
The \code{ggbiplot} package overlays projected observations and variable loadings (arrows) with confidence ellipses:
\vfill
\begin{knitrout}
\definecolor{shadecolor}{rgb}{0.969, 0.969, 0.969}\color{fgcolor}\begin{kframe}
\begin{alltt}
\hlkwd{library}\hldef{(ggbiplot)}
\hldef{p_biplot} \hlkwb{=} \hlkwd{ggbiplot}\hldef{(hockey.pca,} \hlkwc{obs.scale} \hldef{=} \hlnum{1}\hldef{,} \hlkwc{var.scale} \hldef{=} \hlnum{1}\hldef{,}
         \hlkwc{groups} \hldef{= hockey_multi}\hlopt{$}\hldef{position,} \hlkwc{ellipse} \hldef{=} \hlnum{TRUE}\hldef{,}
         \hlkwc{circle} \hldef{=} \hlnum{TRUE}\hldef{,} \hlkwc{alpha} \hldef{=} \hlnum{0.2}\hldef{)} \hlopt{+}
  \hlkwd{theme_minimal}\hldef{()} \hlopt{+}
  \hlkwd{labs}\hldef{(}\hlkwc{title} \hldef{=} \hlsng{"PCA biplot of IIHF players by position"}\hldef{,}
       \hlkwc{color} \hldef{=} \hlsng{"Position"}\hldef{,} \hlkwc{fill} \hldef{=} \hlsng{"Position"}\hldef{)}
\hlkwd{print}\hldef{(p_biplot)}
\end{alltt}
\end{kframe}
\end{knitrout}
\end{frame}

\maxFrameImage{FIGS/U05-plot-pca-ggbiplot-1.pdf}


\begin{frame}\frametitle{Interpreting the biplot}
\textbf{Variable loadings (arrows)}:
\begin{itemize}
\item \code{weight} points along positive PC1 (overall body size)
\item PC2 contrasts \code{age} and \code{bmi} (pointing upward) against \code{height} (pointing downward)
\end{itemize}
\vfill
\textbf{Group ellipses (positions)}:
\begin{itemize}
\item Defencemen tend to have higher PC1 scores (taller and heavier on average)
\item Forwards and goalkeepers cluster at lower PC1 scores
\end{itemize}
\end{frame}


\begin{frame}[fragile]\frametitle{Multivariate analysis with \code{FactoMineR}}
\code{FactoMineR::PCA} supports supplementary qualitative variables (\code{quali.sup = 1} for \code{position}):
\vfill
\begin{knitrout}
\definecolor{shadecolor}{rgb}{0.969, 0.969, 0.969}\color{fgcolor}\begin{kframe}
\begin{alltt}
\hlkwd{library}\hldef{(FactoMineR)}
\hldef{res.pca} \hlkwb{=} \hlkwd{PCA}\hldef{(hockey_multi[,} \hlkwd{c}\hldef{(}\hlsng{"position"}\hldef{,} \hlsng{"height"}\hldef{,} \hlsng{"weight"}\hldef{,}
                               \hlsng{"age"}\hldef{,} \hlsng{"bmi"}\hldef{)],}
              \hlkwc{quali.sup} \hldef{=} \hlnum{1}\hldef{,} \hlkwc{scale.unit} \hldef{=} \hlnum{TRUE}\hldef{,} \hlkwc{graph} \hldef{=} \hlnum{FALSE}\hldef{)}
\end{alltt}
\end{kframe}
\end{knitrout}
\vfill
Provides two key diagnostic visualisations:
\begin{enumerate}
\item \textbf{Variable correlation circle}: correlations with PC1 and PC2 inside the unit circle
\item \textbf{Factor map with group ellipses}: observations and barycentres by position
\end{enumerate}
\end{frame}


\begin{frame}[fragile]\frametitle{Variable correlation circle (\code{FactoMineR})}
\begin{knitrout}
\definecolor{shadecolor}{rgb}{0.969, 0.969, 0.969}\color{fgcolor}\begin{kframe}
\begin{alltt}
\hlkwd{plot.PCA}\hldef{(res.pca,} \hlkwc{axes} \hldef{=} \hlkwd{c}\hldef{(}\hlnum{1}\hldef{,} \hlnum{2}\hldef{),} \hlkwc{choix} \hldef{=} \hlsng{"var"}\hldef{,}
         \hlkwc{title} \hldef{=} \hlsng{"Variable correlation circle (FactoMineR)"}\hldef{)}
\end{alltt}
\end{kframe}
\end{knitrout}
\end{frame}

\maxFrameImage{FIGS/U05-plot-factominer-var-1.pdf}


\begin{frame}[fragile]\frametitle{Position ellipses and barycentres (\code{FactoMineR})}
\begin{knitrout}
\definecolor{shadecolor}{rgb}{0.969, 0.969, 0.969}\color{fgcolor}\begin{kframe}
\begin{alltt}
\hlkwd{plotellipses}\hldef{(res.pca,} \hlkwc{keepvar} \hldef{=} \hlnum{1}\hldef{,} \hlkwc{label} \hldef{=} \hlsng{"quali"}\hldef{,}
             \hlkwc{title} \hldef{=} \hlsng{"Individuals and position barycentres"}\hldef{)}
\end{alltt}
\end{kframe}
\end{knitrout}
\end{frame}

\maxFrameImage{FIGS/U05-plot-factominer-ellipses-1.pdf}




% Save counters for next file

% Save theorem count for next file
\end{document}
